% Lecture notes on the classification of homogeneous plane continua,
% following the formal (Lean 4) proof in this repository.
% Volume 3: General theory of continua.
%
% Compile with:  pdflatex main && pdflatex main   (two runs for the contents and references)
\documentclass[11pt,oneside]{book}

\usepackage[T1]{fontenc}
\usepackage{lmodern}
\usepackage[margin=1.15in]{geometry}
\usepackage{amsmath,amssymb,amsthm}
\usepackage{enumitem}
\usepackage{booktabs}
\usepackage{longtable}
\usepackage{tikz}
\usepackage[hidelinks]{hyperref}

% ---------------------------------------------------------------- theorem styles
\theoremstyle{plain}
\newtheorem{theorem}{Theorem}[chapter]
\newtheorem{lemma}[theorem]{Lemma}
\newtheorem{proposition}[theorem]{Proposition}
\newtheorem{corollary}[theorem]{Corollary}
\newtheorem{claim}[theorem]{Claim}
\theoremstyle{definition}
\newtheorem{definition}[theorem]{Definition}
\newtheorem{example}[theorem]{Example}
\newtheorem{exercise}[theorem]{Exercise}
\theoremstyle{remark}
\newtheorem{remark}[theorem]{Remark}

% ---------------------------------------------------------------- notation
\newcommand{\N}{\mathbb{N}}
\newcommand{\Z}{\mathbb{Z}}
\newcommand{\R}{\mathbb{R}}
\newcommand{\diam}{\operatorname{diam}}
\newcommand{\dist}{\operatorname{dist}}
\newcommand{\interior}{\operatorname{int}}
\newcommand{\cl}[1]{\overline{#1}}
\newcommand{\abs}[1]{\left\lvert #1\right\rvert}
\newcommand{\set}[1]{\left\{ #1 \right\}}
\newcommand{\fl}[1]{\left\lfloor #1\right\rfloor}
% the composant of p in M
\newcommand{\cmp}[2]{\operatorname{cmp}_{#1}(#2)}
% the element of a decomposition containing x
\newcommand{\elt}[1]{[#1]}
% the homeomorphism group
\newcommand{\Homeo}{\mathcal{H}}
% the sup-distance of a homeomorphism from the identity
\newcommand{\mv}{\operatorname{mv}}
% irreducible from A to B
\newcommand{\irr}{\operatorname{irr}}

% Names of declarations in the Lean development (typewriter, underscores allowed).
\newcommand{\lean}[1]{\texttt{\detokenize{#1}}}
% A Lean name ending in a subscript digit (such as IsEssentialSum with subscript 2).
\newcommand{\leansub}[2]{\texttt{\detokenize{#1}}\ensuremath{_{#2}}}
% A paragraph pointing to the formal counterpart of a result.
\newcommand{\formal}[1]{\par\smallskip\noindent{\small\textit{Formal counterpart:} #1}\par\smallskip}
% References to numbered results of Volumes 1 and 2.
\newcommand{\vone}[1]{Volume~1, #1}
\newcommand{\vtwo}[1]{Volume~2, #1}

\setlist{itemsep=2pt,topsep=4pt}

\title{\textbf{Homogeneous Plane Continua}\\[6pt]
\Large Lecture notes following the formal proof\\[18pt]
\large Volume 3\\[4pt]
\Large General theory of continua}
\author{}
\date{}

\begin{document}

\frontmatter
\maketitle
\tableofcontents

\mainmatter
% The introduction is Chapter 0, as in Volumes 1 and 2.
\setcounter{chapter}{-1}
%======================================================================
\chapter{Introduction}\label{ch:intro}
%======================================================================

\section{What this volume is about}

Volumes~1 and~2 dealt with one object, the pseudo-arc. The rest of the classification of
homogeneous plane continua needs a stock of general facts about continua, and this volume
collects them. Most of them are classical and can be found, in some form, in textbooks on
continuum theory; here they are proved in the form in which the formal proof uses them, and
the proofs follow the formal ones.

The volume has seven chapters, which can be read almost independently.

\begin{itemize}
\item \textbf{Chapter~\ref{ch:bumping}: Separated sets and boundary bumping.} Separated sets, the
union of a continuum with one piece of a separation of its complement, nested intersections of
continua, components as intersections of clopen sets, and the \emph{boundary bumping theorem}:
in a continuum $M$, the component of a point $p$ in a proper closed subset $F$ reaches the
boundary of $F$.
\item \textbf{Chapter~\ref{ch:irreducible}: Irreducible continua.} A continuum containing two
points (or meeting two compact sets) contains a subcontinuum that is minimal with this
property; a nondegenerate indecomposable continuum is irreducible between two of its points.
\item \textbf{Chapter~\ref{ch:noncut}: Non-cut points.} R.~L.~Moore's theorem: a nondegenerate
compact connected $T_1$ space has at least two points whose removal leaves it connected, and no
proper connected subset contains all such points.
\item \textbf{Chapter~\ref{ch:triods}: Triods and Sorgenfrey's theorem.} Triods, weak triods and
unicoherence; an indecomposable continuum is not a triod; and Sorgenfrey's theorem: a
nondegenerate unicoherent continuum which is not a triod is irreducible between two of its
points.
\item \textbf{Chapter~\ref{ch:snake}: Snake-like and circle-like continua.} Subcontinua of
snake-like continua are snake-like; snake-like continua are unicoherent and contain no triod;
the $\varepsilon$-map characterization of snake-likeness and its consequence that snake-likeness
is a topological property; circular chains, the circle, and proper subcontinua of circle-like
continua.
\item \textbf{Chapter~\ref{ch:decomp}: Upper semicontinuous decompositions.} Decompositions of a
compact metric space into closed sets, the decomposition space and its metrizability, monotone
decompositions, homeomorphisms permuting the elements, and the theorem of Kuratowski and Fort
that an upper semicontinuous decomposition is continuous at the elements of a residual set.
\item \textbf{Chapter~\ref{ch:effros}: The homeomorphism group and the Effros theorem.} The group
of self-homeomorphisms of a compact metric space is a Polish group; a Polish group acting
transitively on a Baire space acts micro-transitively (the open mapping principle, following
J.~van Mill's elementary proof); and Ungar's corollary for homogeneous compact metric spaces:
close points are moved one to the other by homeomorphisms close to the identity.
\end{itemize}

The main results, in the form used later, are:

\begin{theorem}[Boundary bumping; Theorem~\ref{thm:bump-closed}]
Let $M$ be a continuum, $F\subsetneq M$ a closed proper subset and $p\in F$. Some subcontinuum of
$F$ containing $p$ meets $\cl{M\setminus F}$.
\end{theorem}

\begin{theorem}[Moore; Theorem~\ref{thm:two-noncut}]
A compact connected $T_1$ space with at least two points has two distinct non-cut points.
\end{theorem}

\begin{theorem}[Sorgenfrey; Theorem~\ref{thm:sorgenfrey}]
A nondegenerate unicoherent continuum that is not a triod is irreducible between two of its
points.
\end{theorem}

\begin{theorem}[Theorems~\ref{thm:snake-unicoherent} and~\ref{thm:snake-atriodic}]
A snake-like continuum is hereditarily unicoherent and atriodic.
\end{theorem}

\begin{theorem}[Theorems~\ref{thm:quotient-metrizable} and~\ref{thm:residual}]
The decomposition space of an upper semicontinuous decomposition of a compact metric space into
closed sets is a compact metrizable space, and the decomposition is continuous at the elements of
a residual subset of it.
\end{theorem}

\begin{theorem}[Effros, Ungar; Theorem~\ref{thm:ungar}]
Let $Y$ be a homogeneous compact metric space. For every $\varepsilon>0$ there is $\delta>0$ such
that whenever $d(x,y)<\delta$ there is a homeomorphism $h\colon Y\to Y$ with $h(x)=y$ and
$d(h(z),z)<\varepsilon$ for all $z\in Y$.
\end{theorem}

\section{Where these results are used}

The results of this volume are used throughout Volumes~4--8. The most important uses are:
boundary bumping, in almost every argument about plane continua; irreducibility and Sorgenfrey's
theorem, in the study of snake-like and circle-like continua (Fugate's and Ingram's theorems);
non-cut points, for the end elements of Bing's decomposition of an atriodic continuum; the
fact that indecomposable continua are not triods, together with Moore's theorem on triods in the
plane (a theorem of plane topology), in Hagopian's proof that a homogeneous indecomposable plane continuum is
atriodic; snake-likeness and circle-likeness, in the indecomposable case and for circles of
pseudo-arcs; upper semicontinuous decompositions, for Jones's aposyndetic decomposition of a
decomposable homogeneous continuum (Volume~7); and the Effros theorem, wherever homogeneity is
used quantitatively.

\section{Where the formal proof differs from the literature}\label{sec:intro-diff}

The notes follow the formal proof. It departs from the usual sources in the following places.

\begin{enumerate}
\item \textbf{Non-cut points without a metric.} Moore's theorem is proved, as in Nadler's
\emph{Continuum Theory} (6.3--6.7), for compact connected $T_1$ spaces; the formal statements drop a
hypothesis of nondegeneracy which is implied by the existence of a cut point
(Chapter~\ref{ch:noncut}).
\item \textbf{A false lemma in the proof of Sorgenfrey's theorem.} The lemma that removing a
subcontinuum $H$ from the set $U=X\setminus Z$ of an essential sum leaves a connected set
(Nadler~11.28) is false when $H$ contains $U$: the empty set is not connected. The proof assumes,
without saying so, that $U\setminus H$ is nonempty, and the formal proof adds this hypothesis
(Remark~\ref{rem:1128-false}).
\item \textbf{Sorgenfrey's induction simplified.} In the main step of Sorgenfrey's theorem
(Theorem~\ref{thm:irr-point}) the formal proof absorbs, at stage $n$, the $n$-th point of a dense
sequence, instead of the first point not yet absorbed; only the properties listed in
Lemma~\ref{lem:sorg-step} are used afterwards. Two hypotheses of Nadler's 11.31 and 11.32 are not
needed and are dropped.
\item \textbf{Triods from a core.} In place of Sorgenfrey's theorem that a weak triod contains a
triod, which the formal proof does not use (and does not prove), the later volumes use a simple
``core'' construction (Lemma~\ref{lem:core}).
\item \textbf{Snake-likeness in a compact space.} Snake-likeness is defined with chains of open
sets of the ambient space. Its invariance under homeomorphisms (Theorem~\ref{thm:snake-homeo}) is
proved through $\varepsilon$-maps and the Tietze extension theorem, and needs the target space to
be compact.
\item \textbf{The Effros theorem.} Van Mill's Proposition~2.2 (two analytic fat sets meet) needs
the space to be nonempty; the formal proof proves a relative form (Theorem~\ref{thm:prop22}),
which is what Proposition~3.5 uses. Van Mill's Lemma~3.4 and Proposition~3.5 are proved as
printed, with closures.
\end{enumerate}

\section{Prerequisites and conventions}

The reader should know the basic metric topology of \vone{Chapter~1} (compactness,
connectedness, continua, indecomposability) and, for Chapter~\ref{ch:snake}, the definition of a
chain (\vone{Definition~2.1}). Chapter~\ref{ch:irreducible} uses one result of Volume~2: in an
indecomposable continuum no composant is the whole continuum (\vtwo{Theorem~1.4}).

\begin{itemize}
\item $X$, $S$, $Y$ denote metric spaces unless stated otherwise, $d$ their metrics. $\cl A$ is the
closure of $A$ and $\interior A$ its interior; $B(x,r)$ and $\bar B(x,r)$ are the open and closed
balls.
\item A \emph{continuum} is a nonempty compact connected set (\vone{Definition~1.2}). In the
formal proof this is \lean{HPC.Fold.Q78.IsContinuum}. Most chapters work in a \emph{compact}
metric space $S$; then a set is a continuum if and only if it is nonempty, closed and connected.
\item \emph{Preconnected} means connected or empty, as in \vone{Section~1.3}; the criterion
\vone{(K2)} for closed sets is used constantly: a closed set $C$ is preconnected if and only if
it is not the union of two disjoint nonempty closed sets.
\item A set is \emph{clopen} if it is both closed and open. A subset $W$ of a set $F$ is
\emph{relatively clopen in $F$} if it is clopen in the subspace $F$.
\item $\N=\set{0,1,2,\dots}$ and $[n]=\set{0,\dots,n-1}$.
\end{itemize}

\paragraph{Formal counterparts.} As in Volumes~1 and~2, many results are followed by a line
``Formal counterpart:'' naming the corresponding declaration of the Lean development (all names
begin with \lean{HPC.}). These lines may be skipped. Appendix~\ref{app:formal} collects them.

\paragraph{Exercises.} Each chapter ends with exercises. None of them is used later.

%======================================================================
\chapter{Separated sets and boundary bumping}\label{ch:bumping}
%======================================================================

This chapter collects the elementary tools used in all later chapters: separated sets, unions
of continua, nested intersections of continua, components, and the boundary bumping theorem.
Throughout, $X$ is a metric space.

\section{Separated sets}

\begin{definition}\label{def:separated}
Two sets $A,B\subseteq X$ are \emph{separated} if $\cl A\cap B=\emptyset$ and $A\cap\cl
B=\emptyset$.
\end{definition}

Separated sets are disjoint. Two disjoint closed sets are separated, and so are two disjoint open
sets ($\cl A\cap B=\emptyset$ because $X\setminus B$ is a closed set containing $A$).

\begin{lemma}\label{lem:sep-of-not-conn}
If $C\subseteq X$ is not preconnected, then $C=A\cup B$ with $A,B$ nonempty and separated.
\end{lemma}

\begin{proof}
There are open sets $U,V$ with $C\subseteq U\cup V$, $C\cap U\ne\emptyset$, $C\cap V\neq\emptyset$ and
$C\cap U\cap V=\emptyset$. Put $A=C\cap U$ and $B=C\cap V$; they are nonempty and $C=A\cup B$. If
$x\in\cl A\cap B$, then $x\in V$, and the open neighbourhood $V$ of $x\in\cl A$ contains a point $y$
of $A$; then $y\in C\cap U\cap V$, which is impossible. Symmetrically $A\cap\cl B=\emptyset$.
\end{proof}

\formal{\lean{HPC.Plane.exists_separated_of_not_isPreconnected}.}

\begin{lemma}\label{lem:side}
Let $A,B$ be separated and let $T\subseteq A\cup B$ be preconnected. Then $T\subseteq A$ or
$T\subseteq B$.
\end{lemma}

\begin{proof}
The open sets $U=X\setminus\cl B$ and $V=X\setminus\cl A$ cover $T$: a point of $T$ in $A$ is not in
$\cl B$, and a point of $T$ in $B$ is not in $\cl A$. A point of $T\cap U\cap V$ would lie in $A$ or
in $B$, hence in $\cl A$ or $\cl B$, which is impossible. Since $T$ is preconnected, $T\cap U$ or
$T\cap V$ is empty. If $T\cap V=\emptyset$, every point of $T$ lies in $\cl A$, hence not in $B$
(as $\cl A\cap B=\emptyset$), hence in $A$. The other case is symmetric.
\end{proof}

\formal{\lean{HPC.Plane.subset_side_of_isPreconnected}.}

\begin{lemma}\label{lem:union-continua}
If $A$ and $B$ are continua with $A\cap B\ne\emptyset$, then $A\cup B$ is a continuum.
\end{lemma}

\begin{proof}
$A\cup B$ is compact and nonempty. If $A\cup B$ lies in the union of two disjoint open sets
$U,V$ each meeting $A\cup B$, then $A$ lies in one of them and $B$ in one of them (each is
preconnected); they cannot be the same one, since both meet $A\cup B$; but a point of $A\cap B$
lies in both, which is impossible as $U\cap V\cap(A\cup B)=\emptyset$.
\end{proof}

\formal{\lean{HPC.Plane.isContinuum_union}.}

The next lemma is the basic tool for producing new subcontinua.

\begin{lemma}[A continuum plus one piece]\label{lem:piece}
Let $S$ be a compact metric space, $N\subseteq M$ continua in $S$, and suppose
$M\setminus N=A\cup R$ with $A,R$ separated and $A\neq\emptyset$. Then $N\cup A$ is a continuum.
\end{lemma}

\begin{proof}
\emph{$N\cup A$ is closed.} We have $\cl A\subseteq M$ since $M$ is closed. A point $x\in\cl A$ not in
$N$ lies in $M\setminus N=A\cup R$, and not in $R$ because $\cl A\cap R=\emptyset$; so $x\in A$. Thus
$\cl{N\cup A}=N\cup\cl A\subseteq N\cup A$.

\emph{$N\cup A$ is preconnected.} We use the criterion \vone{(K2)}. Let $N\cup A=F\cup G$ with $F,G$
closed and disjoint. Since $N$ is preconnected, it lies in one of them, say $N\subseteq F$ (the other
case is symmetric). Put $A_G=A\cap G$ and $E=N\cup(A\cap F)\cup R$, so that $M=A_G\cup E$. These two
sets are separated:
\begin{itemize}
\item $\cl{A_G}\subseteq G\cap\cl A$; this misses $R$ (as $\cl A\cap R=\emptyset$) and misses
$N\cup(A\cap F)\subseteq F$ (as $F\cap G=\emptyset$).
\item $\cl E\subseteq F\cup\cl R$, since $N\cup(A\cap F)\subseteq F$ and $F$ is closed; $F$ misses
$A_G\subseteq G$, and $\cl R$ misses $A_G\subseteq A$.
\end{itemize}
$E$ is nonempty (it contains $N$) and $M$ is connected, so $A_G=\emptyset$ by Lemma~\ref{lem:side}
(applied to $T=M$). Hence $A\subseteq F$ and $N\cup A\subseteq F$.

So $N\cup A$ is a nonempty closed preconnected subset of the compact space $S$, that is, a
continuum.
\end{proof}

\formal{\lean{HPC.Plane.isContinuum_union_piece} (Nadler, \emph{Continuum Theory}, 6.3). A
version for an arbitrary connected topological space is Theorem~\ref{thm:nadler63}.}

\section{Nested intersections}

A family $\mathcal C$ of sets is a \emph{chain} (for inclusion) if any two of its members are
comparable: $A\subseteq B$ or $B\subseteq A$.

\begin{lemma}\label{lem:nested}
Let $S$ be a compact metric space and $\mathcal C$ a nonempty chain of continua in $S$. Then
$\bigcap\mathcal C$ is a continuum.
\end{lemma}

\begin{proof}
Put $I=\bigcap\mathcal C$; it is closed. It is nonempty: the members of $\mathcal C$ are compact and
any finitely many of them have a nonempty intersection (the smallest of them), so by compactness
the whole family has a nonempty intersection.

Suppose $I=F\cup G$ with $F,G$ closed, disjoint and nonempty. They are compact, hence at a positive
distance $3\delta$ (\vone{Lemma~1.1}). Let $U$ and $V$ be the open $\delta$-neighbourhoods of $F$
and $G$; they are disjoint. For $C\in\mathcal C$ let $C'=C\setminus(U\cup V)$, a compact set. The
sets $C'$ form a chain, and $\bigcap_{C}C'=I\setminus(U\cup V)=\emptyset$. By compactness some finitely
many of them have empty intersection; the smallest of these, $C_0'$, is then empty. So
$C_0\subseteq U\cup V$. But $C_0\supseteq I$ meets both $U\supseteq F$ and $V\supseteq G$, contradicting
the connectedness of $C_0$.
\end{proof}

\formal{\lean{HPC.T14.isPreconnected_sInter_chain}.}

\section{Components and clopen sets}

Recall (\vtwo{Section~1.4}) that the \emph{component} of a point $p$ in a set $F$ is the largest
preconnected subset of $F$ containing $p$; it is closed when $F$ is.

\begin{lemma}\label{lem:quasi}
Let $K$ be a compact metric space and $p\in K$. The component of $p$ in $K$ is the intersection
$Q$ of all clopen subsets of $K$ containing $p$.
\end{lemma}

\begin{proof}
The component $C$ of $p$ lies in every clopen set $W\ni p$: otherwise $C\cap W$ and $C\setminus W$
would split $C$ into two disjoint nonempty relatively open sets. So $C\subseteq Q$.

For the converse it suffices to show that $Q$ is preconnected (it contains $p$, so it is then
contained in $C$). Let $Q=A\cup B$ with $A,B$ closed and disjoint, $p\in A$; we show $B=\emptyset$.
Choose disjoint open sets $U\supseteq A$, $V\supseteq B$ (for instance the sets of points nearer to
$A$ than to $B$, and conversely). The closed sets $W\setminus(U\cup V)$, $W$ clopen with $p\in W$, have
empty intersection; by compactness finitely many of them do, so $W_0=W_1\cap\dots\cap W_k$
satisfies $W_0\subseteq U\cup V$. The set $W_0\cap U$ is open, and it is closed because it equals
$W_0\setminus V$. It contains $p$, so it is one of the clopen sets defining $Q$, and $Q\subseteq U$.
Then $B\subseteq Q\cap V\subseteq U\cap V=\emptyset$.
\end{proof}

\formal{Mathlib's \lean{connectedComponent_eq_iInter_isClopen}, used inside the proofs below.}

\begin{lemma}\label{lem:clopen-split}
Let $K\subseteq X$ be compact, $p\in K$ and $O$ open, and suppose that every preconnected subset of
$K$ containing $p$ lies in $O$. Then $K=H\cup L$ with $H,L$ compact and disjoint, $p\in H$ and
$H\subseteq O$.
\end{lemma}

\begin{proof}
The component $C$ of $p$ in $K$ lies in $O$. By Lemma~\ref{lem:quasi}, applied in the compact space
$K$, the closed sets $W\setminus O$, $W$ clopen in $K$ with $p\in W$, have empty intersection; by
compactness some finite intersection $H$ of such sets $W$ satisfies $H\subseteq O$. Then $H$ is
clopen in $K$, and $H$ and $L=K\setminus H$ are compact.
\end{proof}

\formal{\lean{HPC.Plane.exists_clopen_split}.}

\begin{theorem}[Separating two closed sets in a compact set]\label{thm:closed-split}
Let $K\subseteq X$ be compact and $P,Q\subseteq K$ closed, and suppose that no preconnected subset of
$K$ meets both $P$ and $Q$. Then $K=H\cup L$ with $H,L$ closed and disjoint, $P\subseteq H$ and
$Q\subseteq L$.
\end{theorem}

\begin{proof}
For each $p\in P$, Lemma~\ref{lem:clopen-split} with $O=X\setminus Q$ gives a splitting
$K=H_p\cup L_p$ into disjoint compact sets with $p\in H_p$ and $H_p\cap Q=\emptyset$. The open sets
$X\setminus L_p$ cover the compact set $P$ (since $p\notin L_p$); choose $p_1,\dots,p_k$ with
$P\subseteq\bigcup_i(X\setminus L_{p_i})$. Put $H=H_{p_1}\cup\dots\cup H_{p_k}$ and $L=K\cap
L_{p_1}\cap\dots\cap L_{p_k}$. Both are closed, and $H\cup L=K$: a point of $K$ in no $H_{p_i}$ lies
in every $L_{p_i}$. They are disjoint because $H_{p_i}\cap L_{p_i}=\emptyset$. A point of $P$ lies
outside some $L_{p_i}$, hence in $H_{p_i}\subseteq H$. Finally $Q\cap H=\emptyset$, so $Q\subseteq L$.
\end{proof}

\formal{\lean{HPC.Plane.exists_closed_split_of_no_connected} (Moore's Theorem~44 in
\emph{Foundations of Point Set Theory}). A variant with a continuum in place of the preconnected
set, in a compact space, is \lean{HPC.Plane.closed_dichotomy}.}

\section{Boundary bumping}

We use the clopen form of boundary bumping proved in Volume~2:

\begin{quote}
\vtwo{Lemma~1.5}. \emph{Let $M$ be preconnected, $L\subseteq M$ with $L\ne M$, and let $W\subseteq L$ be
nonempty, closed, and relatively open in $L$. Then $W$ meets $\cl{M\setminus L}$.}
\end{quote}

\begin{theorem}[Boundary bumping]\label{thm:bump-closed}
Let $M\subseteq X$ be preconnected, $F\subseteq M$ compact with $F\ne M$, and $p\in F$. Then there is a
continuum $K$ with $p\in K\subseteq F$ and $K\cap\cl{M\setminus F}\ne\emptyset$. One can take for
$K$ the component of $p$ in $F$.
\end{theorem}

\begin{proof}
Let $K$ be the component of $p$ in $F$; it is closed in the compact set $F$, so it is a continuum.
Suppose $K\cap\cl{M\setminus F}=\emptyset$. The set $H=F\cap\cl{M\setminus F}$ is compact and misses
$K$. By Lemma~\ref{lem:quasi} (in the compact space $F$), the closed sets $H\cap W$, $W$ clopen in $F$
with $p\in W$, have empty intersection; by compactness there is a finite intersection $W$ of such
sets with $H\cap W=\emptyset$. This $W$ is clopen in $F$: it is closed in $X$ (closed in the compact
$F$), and $W=F\cap O$ for some open $O\subseteq X$. It is nonempty ($p\in W$). By \vtwo{Lemma~1.5}
(with $L=F$) $W$ meets $\cl{M\setminus F}$; since $W\subseteq F$, it meets $H$, a contradiction.
\end{proof}

\formal{\lean{HPC.Jones.exists_continuum_bump_compact}; for $M$ a continuum in a compact space and
$F$ closed, \lean{HPC.Plane.exists_continuum_bump_closed} (Nadler 5.4).}

\begin{corollary}[Boundary bumping, open form]\label{cor:bump-open}
Let $S$ be a compact metric space, $L$ a continuum, $O$ an open set, $p\in L\cap O$ and
$L\not\subseteq O$. Then some subcontinuum $K$ of $L$ with $p\in K\subseteq\cl{L\cap O}$ meets
$L\setminus O$.
\end{corollary}

\begin{proof}
Put $F=\cl{L\cap O}\subseteq L$. If $F=L$, take $K=L$. Otherwise Theorem~\ref{thm:bump-closed} gives a
continuum $K\ni p$ inside $F$ meeting $\cl{L\setminus F}$. Now $L\setminus F\subseteq L\setminus O$
(a point of $L\cap O$ lies in $F$), and $L\setminus O$ is closed, so $\cl{L\setminus F}\subseteq
L\setminus O$.
\end{proof}

\formal{\lean{HPC.Plane.exists_continuum_bump_open} (Nadler 5.7).}

\begin{corollary}[Boundary bumping, metric form]\label{cor:bump-ball}
Let $M$ be a continuum, $p,x\in M$, $s>0$ and $d(p,x)\ge s$. Some subcontinuum $K$ of
$M\setminus B(x,s)$ contains $p$ and meets $\bar B(x,s)$.
\end{corollary}

\begin{proof}
Apply Theorem~\ref{thm:bump-closed} to $F=M\setminus B(x,s)$, which is compact, contains $p$, and
misses $x\in M$. The closure of $M\setminus F=M\cap B(x,s)$ lies in $\bar B(x,s)$.
\end{proof}

\formal{\lean{HPC.Plane.exists_subcontinuum_meeting_closedBall}.}

\begin{example}
Let $M=[0,1]$, $F=[0,\tfrac12]\cup\set1$ and $p=0$. The component of $0$ in $F$ is $[0,\frac12]$, which
meets $\cl{M\setminus F}=[\frac12,1]$ at $\frac12$. The component of $1$ in $F$ is $\set1$, which
also meets $[\frac12,1]$: boundary bumping holds for every point of $F$.
\end{example}

\section*{Exercises}

\begin{exercise}
Show that Lemma~\ref{lem:nested} fails without compactness: find a decreasing sequence of closed
connected subsets of $\R^2$ whose intersection is not connected.
\end{exercise}

\begin{exercise}
Let $M$ be a continuum and $U\subseteq M$ a nonempty open subset of $M$ with $U\neq M$. Show that every
component of $\cl U$ meets the boundary of $U$ in $M$.
\end{exercise}

\begin{exercise}
Deduce from Theorem~\ref{thm:bump-closed} that a nondegenerate continuum contains no isolated
points, and in fact that each point of it lies in arbitrarily small nondegenerate subcontinua.
\end{exercise}

\begin{exercise}
Show that Lemma~\ref{lem:quasi} fails for the non-compact space $\set{(0,0),(0,1)}\cup\bigcup_n
\set{1/n}\times[0,1]\subseteq\R^2$ with $p=(0,0)$.
\end{exercise}

%======================================================================
\chapter{Irreducible continua}\label{ch:irreducible}
%======================================================================

Throughout this chapter $S$ is a compact metric space.

\section{Definitions and existence}

\begin{definition}\label{def:irreducible}
Let $W$ be a continuum and $a,b$ points. $W$ is \emph{irreducible between $a$ and $b$} if $a,b\in W$
and no subcontinuum of $W$ other than $W$ contains both $a$ and $b$. A continuum is
\emph{irreducible} if it is irreducible between two of its points.
\end{definition}

\formal{\lean{HPC.IsIrreducibleBetween} (in \lean{HPC/PseudoArc/Theorem14.lean}).}

\begin{definition}\label{def:irr-sets}
Let $C$ be a continuum and $A,B$ sets. $C$ is \emph{irreducible from $A$ to $B$}, written
$C=\irr(A,B)$, if $C$ meets $A$ and $B$, and no subcontinuum of $C$ other than $C$ meets both $A$
and $B$.
\end{definition}

\formal{\lean{HPC.Plane.IrrFromTo} (Nadler 11.29).}

For points, $C$ is irreducible between $a$ and $b$ exactly when $C=\irr(\set a,\set b)$.

\begin{example}\label{ex:arc-irr}
The arc $[0,1]$ is irreducible between $0$ and $1$: a connected subset of $\R$ containing $0$ and
$1$ contains $[0,1]$. It is not irreducible between $0$ and $\frac12$, since $[0,\frac12]$ is a
smaller subcontinuum containing both. A circle is not irreducible: any two of its points lie in
a closed half-circle, which is a proper subcontinuum.
\end{example}

We use Zorn's lemma in the following form: \emph{if $\mathcal S$ is a nonempty family of sets such
that the intersection of every nonempty chain in $\mathcal S$ contains a member of $\mathcal S$,
then $\mathcal S$ has a minimal member} (one containing no other member of $\mathcal S$).

\begin{lemma}\label{lem:nested-meets}
Let $\mathcal C$ be a nonempty chain of continua in $S$ and $D\subseteq S$ compact. If every member
of $\mathcal C$ meets $D$, then $\bigcap\mathcal C$ meets $D$.
\end{lemma}

\begin{proof}
The compact sets $C\cap D$, $C\in\mathcal C$, form a chain of nonempty sets; any finitely many of
them have nonempty intersection (the smallest one). By compactness their whole intersection,
$\bigcap\mathcal C\cap D$, is nonempty.
\end{proof}

\begin{theorem}[Existence of irreducible subcontinua]\label{thm:irr-exists}
Let $X_0$ be a continuum in $S$ and $A,B$ compact sets meeting $X_0$. Then $X_0$ contains a
subcontinuum $C=\irr(A,B)$. In particular, a continuum containing two points $a,b$ contains a
subcontinuum irreducible between $a$ and $b$.
\end{theorem}

\begin{proof}
Let $\mathcal S$ be the family of subcontinua of $X_0$ meeting both $A$ and $B$; it contains $X_0$.
Let $\mathcal C\subseteq\mathcal S$ be a nonempty chain. Its intersection is a continuum
(Lemma~\ref{lem:nested}), lies in $X_0$, and meets $A$ and $B$ (Lemma~\ref{lem:nested-meets}); so
it is a member of $\mathcal S$. By Zorn's lemma $\mathcal S$ has a minimal member $C$. If $K\subseteq
C$ is a subcontinuum meeting $A$ and $B$, then $K\in\mathcal S$, so $K=C$ by minimality. Thus
$C=\irr(A,B)$. For two points take $A=\set a$, $B=\set b$.
\end{proof}

\formal{\lean{HPC.Plane.exists_irrFromTo} (Nadler 11.30) and \lean{HPC.exists_irreducible}.}

\section{Indecomposable continua are irreducible}

\begin{theorem}\label{thm:indec-irr}
Let $X_0$ be a nondegenerate indecomposable continuum in $S$ and $p\in X_0$. There is $q\in X_0$,
$q\neq p$, such that $X_0$ is irreducible between $p$ and $q$. In fact any $q$ outside the composant
of $p$ will do.
\end{theorem}

\begin{proof}
By \vtwo{Theorem~1.4} there is $q\in X_0$ lying in no proper subcontinuum of $X_0$ containing $p$.
If $K\subseteq X_0$ is a subcontinuum containing $p$ and $q$, then $K$ is not proper, that is,
$K=X_0$. So $X_0$ is irreducible between $p$ and $q$. If $q=p$, the proper subcontinuum $\set p$
(proper because $X_0$ is nondegenerate) would contain $p$ and $q$; so $q\neq p$.
\end{proof}

\formal{\lean{HPC.Plane.exists_irreducible_of_indecomposable} (the case of Nadler 11.15 used
later).}

\begin{remark}
In an indecomposable continuum every composant is dense, and there are uncountably many
composants; so an indecomposable continuum is irreducible between ``most'' pairs of its points.
We shall not need this.
\end{remark}

\begin{proposition}\label{prop:irr-homeo}
Let $M$ be a continuum, $T\colon M\to M$ a homeomorphism, and $W\subseteq M$ a subcontinuum
irreducible between $a$ and $b$. Then $T(W)$ is irreducible between $T(a)$ and $T(b)$.
\end{proposition}

\begin{proof}
$T(W)$ is a continuum containing $T(a)$ and $T(b)$. If $K\subseteq T(W)$ is a subcontinuum containing
them, then $T^{-1}(K)\subseteq W$ is a subcontinuum containing $a$ and $b$, so $T^{-1}(K)=W$ and
$K=T(W)$.
\end{proof}

\formal{\lean{HPC.IsIrreducibleBetween.homImage}.}

\section*{Exercises}

\begin{exercise}
Show that the $\sin(1/x)$-continuum, the closure in $\R^2$ of $\set{(x,\sin(1/x)):0<x\le1}$, is
irreducible between $(1,\sin1)$ and any point of the limit segment $\set0\times[-1,1]$, but is not
irreducible between $(0,0)$ and $(0,1)$.
\end{exercise}

\begin{exercise}
Show that a continuum irreducible between $a$ and $b$ is not the union of three subcontinua with a
common point, none of which is contained in the union of the other two. (This is
Lemma~\ref{lem:weak-not-irr} below.)
\end{exercise}

\begin{exercise}
Give an example of a continuum $X_0$ and closed sets $A,B$ for which $\irr(A,B)$ is not unique.
\end{exercise}

%======================================================================
\chapter{Non-cut points}\label{ch:noncut}
%======================================================================

In this chapter $X$ is a topological space, not necessarily metric. The main theorem is due to
R.~L.~Moore; we follow Nadler (\emph{Continuum Theory}, 6.3--6.7), as the formal proof does. The
space will be compact, connected and $T_1$ (points are closed). It is applied later not to a
continuum but to the decomposition space of a decomposition of a continuum
(Remark~\ref{rem:noncut-use}), which is why we do not assume a metric.

\section{Separations and cut points}

\begin{definition}\label{def:separation}
Let $Y\subseteq X$. A \emph{separation} $Y=P\mid Q$ is a pair of nonempty disjoint sets $P,Q$ with
$P\cup Q=Y$, both open in the subspace $Y$. A point $x\in X$ is a \emph{cut point} of $X$ if
$X\setminus\set x$ is not connected; otherwise it is a \emph{non-cut point}.
\end{definition}

\formal{\lean{HPC.Plane.IsSeparation}, \lean{HPC.Plane.IsCutPoint} (in
\lean{HPC/Plane/NonCut51Gen.lean}).}

\begin{lemma}\label{lem:sep-basic}
Let $Y=P\mid Q$ be a separation.
\begin{enumerate}[label=(\alph*)]
\item $P$ and $Q$ are separated: $\cl P\cap Q=\emptyset=P\cap\cl Q$.
\item If $Y$ is open in $X$, then $P$ and $Q$ are open in $X$.
\item A preconnected subset of $Y$ lies in $P$ or in $Q$.
\item If $Y$ is nonempty and not connected and $c\in Y$, then $Y$ has a separation $Y=P\mid Q$ with
$c\in Q$.
\end{enumerate}
\end{lemma}

\begin{proof}
(a) Write $Q=Y\cap t$ with $t$ open in $X$. Then $t\cap P=\emptyset$, so $t\cap\cl P=\emptyset$, so
$Q\cap\cl P=\emptyset$. Symmetrically for the other half.
(b) $Q=Y\cap t$ is an intersection of two open sets.
(c) is Lemma~\ref{lem:side} (whose proof uses only the topology), with (a).
(d) As $Y$ is not preconnected there are open $u,v\subseteq X$ with $Y\subseteq u\cup v$, $Y\cap u\cap
v=\emptyset$ and both $Y\cap u$, $Y\cap v$ nonempty. These form a separation; name them so that
$c$ lies in the second.
\end{proof}

\formal{\lean{HPC.Plane.IsSeparation.disjoint_closure_left},
\lean{HPC.Plane.IsSeparation.isOpen_right}, \lean{HPC.Plane.IsSeparation.subset_or_subset},
\lean{HPC.Plane.exists_separation_of_not_isConnected}.}

\begin{theorem}\label{thm:nadler63}
Let $X$ be connected, $C\subseteq X$ connected, and $X\setminus C=A\mid B$. Then $A\cup C$ is
connected.
\end{theorem}

\begin{proof}
$A\cup C$ is nonempty. Let $u,v$ be open sets with $A\cup C\subseteq u\cup v$ and $(A\cup C)\cap u\cap
v=\emptyset$; we must show $A\cup C\subseteq u$ or $A\cup C\subseteq v$. Since $C$ is connected, it lies
in one of them, say $C\subseteq u$. Put
\[
P=A\cap v,\qquad W=(A\cap u)\cup C .
\]
Then $W\subseteq u$, and $X=P\cup W\cup B$ (a point not in $C$ lies in $A$ or $B$, and a point of $A$ lies
in $u$ or $v$). The sets $P$ and $W\cup B$ are separated:
\begin{itemize}
\item $\cl P\cap W=\emptyset$: the open set $u$ misses $P$ (a point of $P\cap u$ would be in $(A\cup
C)\cap u\cap v$), hence misses $\cl P$, and $W\subseteq u$.
\item $\cl P\cap B\subseteq\cl A\cap B=\emptyset$ by Lemma~\ref{lem:sep-basic}(a).
\item $P\cap\cl W=\emptyset$: the open set $v$ misses $W$ (as $W\subseteq(A\cup C)\cap u$), hence misses
$\cl W$, and $P\subseteq v$.
\item $P\cap\cl B\subseteq A\cap\cl B=\emptyset$.
\end{itemize}
$B$ is nonempty and $X$ is connected, so $P=\emptyset$ by Lemma~\ref{lem:side}. Hence $A\subseteq u$ and
$A\cup C\subseteq u$.
\end{proof}

\formal{\lean{HPC.Plane.isConnected_union_of_separation} (Nadler 6.3).}

\begin{lemma}\label{lem:nadler64}
Let $X$ be connected, $X\setminus\set x=K\mid L$ and $X\setminus\set y=M\mid N$, with $x\in M$ and
$y\in K$. Then $N\cup\set y\subseteq K$.
\end{lemma}

\begin{proof}
$x\ne y$ since $x\in M\subseteq X\setminus\set y$. By Theorem~\ref{thm:nadler63} (with $C=\set y$ and the
separation $X\setminus\set y=N\mid M$) the set $N\cup\set y$ is connected. It lies in $X\setminus\set
x$: $x\in M$ is not in $N$, and $x\neq y$. By Lemma~\ref{lem:sep-basic}(c) it lies in $K$ or in $L$; it
contains $y\in K$, so it lies in $K$.
\end{proof}

\formal{\lean{HPC.Plane.nadler_6_4}.}

\section{Moore's theorem}

\begin{theorem}[Moore]\label{thm:noncut-side}
Let $X$ be a compact connected $T_1$ space, $c\in X$, and $X\setminus\set c=U\mid V$. Then $U$
contains a non-cut point of $X$, and so does $V$.
\end{theorem}

\begin{proof}
By symmetry it suffices to treat $U$. Suppose every point of $U$ is a cut point. For each $x\in U$,
Lemma~\ref{lem:sep-basic}(d) gives a separation $X\setminus\set x=U_x\mid V_x$ with $c\in V_x$ (note
$c\neq x$). Put $A(x)=U_x\cup\set x$. We prove four facts.

\emph{(i) $A(x)\subseteq U$.} Apply Lemma~\ref{lem:nadler64} to the separations $X\setminus\set
c=U\mid V$ and $X\setminus\set x=V_x\mid U_x$: we have $c\in V_x$ and $x\in U$, so $U_x\cup\set x\subseteq
U$.

\emph{(ii) $A(x)$ is closed.} Its complement is $V_x$, which is open in $X$ by
Lemma~\ref{lem:sep-basic}(b), since $X\setminus\set x$ is open ($X$ is $T_1$).

\emph{(iii) If $y\in U_x$, then $A(y)\subseteq U_x$; in particular $x\notin A(y)$.} By (i), $y\in U$. Since
$x\neq y$, $x$ lies in $U_y$ or in $V_y$. If $x\in U_y$, Lemma~\ref{lem:nadler64} applied to
$X\setminus\set y=U_y\mid V_y$ and $X\setminus\set x=U_x\mid V_x$ (with $y\in U_x$ and $x\in U_y$) gives
$V_x\cup\set x\subseteq U_y$; but $c\in V_x$ and $c\in V_y$, a contradiction. So $x\in V_y$, and
Lemma~\ref{lem:nadler64} applied to $X\setminus\set x=U_x\mid V_x$ and $X\setminus\set y=V_y\mid U_y$
(with $x\in V_y$ and $y\in U_x$) gives $U_y\cup\set y\subseteq U_x$. Finally $x\notin U_x$.

\emph{(iv) The family $\set{A(x):x\in U}$ has a minimal member.} We use Zorn's lemma in the form stated before
Lemma~\ref{lem:nested-meets}. The family is nonempty ($U\neq\emptyset$). Let $\mathcal C$ be a
nonempty chain in it. Its members are closed and nonempty, any finitely many have nonempty
intersection, and $X$ is compact; so there is a point $p$ in all of them. Then $p\in U$ by (i). For
each member $A(x)$ of $\mathcal C$, either $p\in U_x$, and then $A(p)\subseteq U_x\subseteq A(x)$ by (iii),
or $p=x$ and $A(p)=A(x)$. So $A(p)$ is a member of the family contained in $\bigcap\mathcal C$.

Now let $A(x)$ be minimal, and take $q\in U_x$ (it is nonempty). By (iii) $A(q)\subseteq U_x\subseteq
A(x)$, so $A(q)=A(x)$ by minimality, and $x\in A(x)=A(q)$, contradicting (iii).
\end{proof}

\formal{\lean{HPC.Plane.exists_nonCutPoint_left} and \lean{HPC.Plane.exists_nonCutPoint_of_cut}
(Nadler 6.6). Nadler assumes $X$ nondegenerate; this follows from the separation and is not
needed.}

\begin{theorem}\label{thm:two-noncut}
A compact connected $T_1$ space with at least two points has two distinct non-cut points.
\end{theorem}

\begin{proof}
If $X$ has no cut point, any two distinct points will do. If $c$ is a cut point, take $d\ne c$ and a
separation $X\setminus\set c=U\mid V$ (Lemma~\ref{lem:sep-basic}(d)); Theorem~\ref{thm:noncut-side}
gives non-cut points $u\in U$ and $v\in V$, which are distinct as $U\cap V=\emptyset$.
\end{proof}

\formal{\lean{HPC.Plane.exists_two_nonCutPoints}.}

\begin{theorem}\label{thm:noncut-not-in}
Let $X$ be a compact connected $T_1$ space and $Z\subsetneq X$ a connected proper subset. Then some
non-cut point of $X$ lies outside $Z$.
\end{theorem}

\begin{proof}
Suppose all non-cut points lie in $Z$. Take $c\notin Z$; it is a cut point. Take $z\in Z$ and a
separation $X\setminus\set c=U\mid V$. The connected set $Z\subseteq X\setminus\set c$ lies in $U$ or in
$V$, say in $U$. By Theorem~\ref{thm:noncut-side} $V$ contains a non-cut point, which is then in
$Z\subseteq U$; but $U\cap V=\emptyset$.
\end{proof}

\formal{\lean{HPC.Plane.not_subset_of_nonCutPoints} (Nadler 6.7).}

\begin{example}
The non-cut points of $[0,1]$ are $0$ and $1$. Every point of a circle is a non-cut point. In the
letter T (three segments with a common end point $o$), the non-cut points are the three free end
points; $o$ is a cut point whose complement has three components.
\end{example}

\begin{remark}\label{rem:noncut-use}
In the later volumes (Bing's theorem on atriodic hereditarily decomposable continua) these
results are applied to the decomposition space of an upper semicontinuous decomposition of a
continuum into subcontinua (Chapter~\ref{ch:decomp}): the decomposition space is compact,
connected and $T_1$, and its non-cut points correspond to the ``end elements'' of the
decomposition (\lean{HPC.Plane.exists_two_endElts}).
\end{remark}

\section*{Exercises}

\begin{exercise}
Show that compactness cannot be dropped from Theorem~\ref{thm:two-noncut}: every point of $\R$ is
a cut point. Where does the proof of Theorem~\ref{thm:noncut-side} use compactness, and where
does it use the $T_1$ property?
\end{exercise}

\begin{exercise}
Show that a nondegenerate continuum with exactly two non-cut points is an arc. (Hard; this is a
classical characterization of the arc. It is not used in these notes.)
\end{exercise}

\begin{exercise}
Let $X$ be a compact connected $T_1$ space and $c$ a cut point. Show that each component of
$X\setminus\set c$ contains a non-cut point of $X$, if $X\setminus\set c$ has finitely many components.
\end{exercise}

%======================================================================
\chapter{Triods and Sorgenfrey's theorem}\label{ch:triods}
%======================================================================

Throughout this chapter $S$ is a compact metric space and all continua lie in $S$. The chapter
follows Nadler, \emph{Continuum Theory}, 11.22--11.35, as the formal proof does, with the
corrections listed in Section~\ref{sec:intro-diff}.

\section{Triods, unicoherence and essential sums}

\begin{definition}\label{def:triod}
A \emph{triod} is a continuum $T$ containing a subcontinuum $N$ (the \emph{core}) such that
$T\setminus N=A\cup B\cup C$ with $A,B,C$ nonempty and pairwise separated
(Definition~\ref{def:separated}). A continuum is \emph{atriodic} if no subcontinuum of it is a
triod.
\end{definition}

\formal{\lean{HPC.Plane.IsTriod}, \lean{HPC.Plane.IsAtriodic} (in
\lean{HPC/Plane/Atriodic51Defs.lean}); for a triod in a non-compact space,
\lean{HPC.Hagopian.IsTriodSet}.}

The letter T is a triod (with core its centre point). An arc and a circle are atriodic; this will
follow from Chapter~\ref{ch:snake} for the arc.

\begin{definition}\label{def:unicoherent}
A continuum $M$ is \emph{unicoherent} if whenever $M=H\cup K$ with $H,K$ subcontinua, $H\cap K$ is
connected. It is \emph{hereditarily unicoherent} if every nondegenerate subcontinuum of $M$ is
unicoherent.
\end{definition}

\formal{\lean{HPC.Plane.IsUnicoherent}, \lean{HPC.Plane.IsHereditarilyUnicoherent}.}

An arc is unicoherent; a circle is not (two closed half-circles meet in two points).

\begin{definition}\label{def:esum}
Let $X,Y,Z,X_1,X_2,X_3$ be continua.
\begin{enumerate}[label=(\alph*)]
\item $X=Y\oplus Z$ (an \emph{essential sum}) means $X=Y\cup Z$, $Y\ne X$ and $Z\neq X$. So $X$ is
decomposable if and only if it is an essential sum of two subcontinua.
\item $X=X_1\oplus X_2\oplus X_3$ means $X=X_1\cup X_2\cup X_3$ and no $X_i$ is contained in the union
of the other two.
\item $X$ is a \emph{weak triod} if $X=X_1\oplus X_2\oplus X_3$ with $X_1\cap X_2\cap X_3\ne\emptyset$.
\end{enumerate}
\end{definition}

\formal{\leansub{HPC.Plane.IsEssentialSum}{2}, \leansub{HPC.Plane.IsEssentialSum}{3},
\lean{HPC.Plane.IsWeakTriod} (in \lean{HPC/Plane/SorgenfreyDefs.lean}); the constructor used in
all proofs below is \lean{HPC.Plane.isWeakTriod_of}.}

To prove that $X$ is a weak triod it suffices to give three continua $X_1,X_2,X_3$ with union $X$,
a common point, and for each $i$ a point of $X_i$ outside the other two. The union is then
automatically a continuum (Lemma~\ref{lem:union-continua}). This is how weak triods are produced in
every proof of this chapter.

\paragraph{Notation for an essential sum.} When $X=Y\oplus Z$ we write
\[
U=X\setminus Z,\qquad V=X\setminus Y .
\]

\begin{lemma}\label{lem:esum-basic}
Let $X=Y\oplus Z$. Then:
\begin{enumerate}[label=(\alph*)]
\item $U$ and $V$ are nonempty, $U\subseteq Y$, $V\subseteq Z$, $\cl U\subseteq Y$ and $\cl V\subseteq Z$;
\item $X=\cl U\cup Z$, $\cl U\ne X$, and $\cl U$ misses $Z\setminus Y=V$;
\item $\cl U\cap Z\ne\emptyset$.
\end{enumerate}
\end{lemma}

\begin{proof}
(a) If $U=\emptyset$ then $X\subseteq Z$, so $Z=X$; similarly for $V$. A point of $X$ not in $Z$ is in
$Y$, so $U\subseteq Y$, and $\cl U\subseteq Y$ as $Y$ is closed. Symmetrically for $V$.
(b) A point of $X$ outside $Z$ is in $U$. Since $\cl U\subseteq Y$, $\cl U$ misses $X\setminus Y=V$, which
is nonempty and contained in $X$; so $\cl U\neq X$.
(c) $X=\cl U\cup Z$ with both sets closed and nonempty, and $X$ is connected.
\end{proof}

\formal{The lemmas in the namespace \leansub{HPC.Plane.IsEssentialSum}{2}
(\lean{diff_nonempty}, \lean{closure_diff_subset_left}, \lean{eq_closure_union},
\lean{closure_diff_ne}, \lean{closure_diff_disjoint}, \lean{closure_inter_nonempty}).}

\section{Triods and weak triods}

\begin{lemma}[A triod with a given core]\label{lem:core}
Let $N,X_1,X_2,X_3$ be continua such that each $X_i$ meets $N$, $X_i\cap X_j\subseteq N$ for $i\ne j$,
and no $X_i$ is contained in $N$. Then $N\cup X_1\cup X_2\cup X_3$ is a triod with core $N$.
\end{lemma}

\begin{proof}
The union is a continuum by Lemma~\ref{lem:union-continua}. Its complement of $N$ is the union of
the three nonempty sets $X_i\setminus N$. Since $\cl{X_i\setminus N}\subseteq X_i$, a point of
$\cl{X_i\setminus N}\cap(X_j\setminus N)$ would lie in $X_i\cap X_j\subseteq N$; so the three sets are
pairwise separated.
\end{proof}

\formal{\lean{HPC.Plane.isTriod_of_core}.}

\begin{proposition}\label{prop:triod-weak}
Every triod is a weak triod.
\end{proposition}

\begin{proof}
Let $T\setminus N=A\cup B\cup C$ be as in Definition~\ref{def:triod}. By Lemma~\ref{lem:piece} (with
$M=T$ and $R=B\cup C$, which is separated from $A$) the set $X_1=N\cup A$ is a continuum; likewise
$X_2=N\cup B$ and $X_3=N\cup C$. Their union is $T$, and they share the points of $N$. A point $a\in A$ is
not in $N$, and not in $B\cup C$ because $A$ is separated from $B$ and $C$; so $a\notin X_2\cup X_3$.
Similarly for $B$ and $C$.
\end{proof}

\formal{\lean{HPC.Plane.isWeakTriod_of_isTriod} (Nadler 11.25).}

\begin{corollary}\label{cor:indec-not-triod}
An indecomposable continuum is not a triod.
\end{corollary}

\begin{proof}
If $X=X_1\oplus X_2\oplus X_3$ with a common point $c$, then $X_2\cup X_3$ is a continuum
(Lemma~\ref{lem:union-continua}) and $X=X_1\cup(X_2\cup X_3)$. Here $X_1\ne X$, since otherwise
$X_2\subseteq X_1\cup X_3$; and $X_2\cup X_3\ne X$, since otherwise $X_1\subseteq X_2\cup X_3$. So $X$ is
decomposable. Apply Proposition~\ref{prop:triod-weak}.
\end{proof}

\formal{\lean{HPC.Hagopian.not_isTriod_univ_of_indecomposable} (for the whole space, which is all
that is used).}

\begin{lemma}\label{lem:weak-not-irr}
A weak triod is not irreducible between any two points.
\end{lemma}

\begin{proof}
Let $X=X_1\oplus X_2\oplus X_3$ with a common point $c$, and $p,q\in X$. Each of $p,q$ lies in some
$X_i$, so there are $i\neq j$ with $p,q\in X_i\cup X_j$. This is a continuum (it contains $c$) and is
not $X$: with $k$ the third index, $X_k\not\subseteq X_i\cup X_j$.
\end{proof}

\formal{\lean{HPC.Plane.not_irreducible_of_isWeakTriod}.}

\begin{proposition}\label{prop:uni-triod}
A unicoherent continuum is a triod if and only if it is a weak triod.
\end{proposition}

\begin{proof}
One direction is Proposition~\ref{prop:triod-weak}. Conversely let $X=X_1\oplus X_2\oplus X_3$ be
unicoherent, with a common point $c$. For $\set{i,j,k}=\set{1,2,3}$ the set $X_j\cup X_k$ is a
continuum, and $X=X_i\cup(X_j\cup X_k)$, so by unicoherence $Z_i=X_i\cap(X_j\cup X_k)$ is a continuum.
The $Z_i$ contain $c$, so $N=Z_1\cup Z_2\cup Z_3$ is a continuum. Let $A_i=X\setminus(X_j\cup X_k)$,
the set of points lying in $X_i$ only; it is nonempty because $X_i\not\subseteq X_j\cup X_k$.

A point of $X$ outside $N$ lies in some $X_i$ and in no other $X_j$ (otherwise it would be in $Z_i$),
so $X\setminus N=A_1\cup A_2\cup A_3$; and the $A_i$ miss $N$. Finally $\cl{A_i}\subseteq X_i$, and $A_j$
misses $X_i$ for $j\ne i$, so the $A_i$ are pairwise separated. Thus $X$ is a triod with core $N$.
\end{proof}

\formal{\lean{HPC.Plane.isTriod_iff_isWeakTriod} (Nadler 11.26).}

\begin{remark}
Without unicoherence the converse of Proposition~\ref{prop:triod-weak} fails: the union of a circle
and an arc meeting it in one end point $c$ is a weak triod (the arc and two arcs of the circle
through $c$), but it is not a triod. Sorgenfrey proved that every weak triod \emph{contains} a
triod. The formal proof does not use this theorem and does not prove it; the
later volumes produce triods by Lemma~\ref{lem:core} instead.
\end{remark}

\section{Connected differences}

\begin{theorem}\label{thm:diff-connected}
Let $X=Y\oplus Z$, and suppose $X$ is not a weak triod. Then $U=X\setminus Z$ and $V=X\setminus Y$ are
connected. In particular $\cl U$ and $\cl V$ are continua.
\end{theorem}

\begin{proof}
By symmetry we prove that $V$ is connected; it is nonempty by Lemma~\ref{lem:esum-basic}.

\emph{Step 1. If $V=P\cup Q$ with $P,Q$ separated and $Q\neq\emptyset$, then $P$ is preconnected.}
Otherwise $P=W_1\cup W_2$ with $W_1,W_2$ nonempty and separated (Lemma~\ref{lem:sep-of-not-conn}). The
three sets $W_1$, $W_2$, $Q$ are pairwise separated (for instance $\cl{W_1}\cap Q\subseteq\cl P\cap
Q=\emptyset$), and $X\setminus Y=W_1\cup W_2\cup Q$. By Lemma~\ref{lem:piece} the sets $Y\cup W_1$,
$Y\cup W_2$ and $Y\cup Q$ are continua. Their union is $X$, they share the points of $Y$, and a point
of $W_1$ lies neither in $Y\cup W_2$ nor in $Y\cup Q$; similarly for $W_2$ and $Q$. So $X$ is a weak
triod, which is excluded.

\emph{Step 2. Suppose $V$ is not preconnected}, and write $V=E\cup F$ with $E,F$ nonempty and
separated. By Step~1 both $E$ and $F$ are preconnected, so $\cl E$ and $\cl F$ are continua, contained
in $Z$ (since $V\subseteq Z$). We claim: \emph{no preconnected subset $K'$ of $Y\cap Z$ meets both
$\cl E$ and $\cl F$.} Suppose $K'$ does, and let $K=\cl{K'}\subseteq Y\cap Z$, a continuum. Then $K\cup\cl
E$, $K\cup\cl F$ and $Y$ are continua with union $X$ (it contains $Y\cup E\cup F=Y\cup V$), with a common
point (any point of $K'$). A point $a\in E$ is not in $Y$, hence not in $K$, and not in $\cl F$; so
$a\notin(K\cup\cl F)\cup Y$. Similarly a point of $F$ is outside the other two. A point $y\in U$ lies in
$Y$, but not in $Z\supseteq K\cup\cl E\cup\cl F$. So $X$ is a weak triod, which is excluded.

By the claim and Theorem~\ref{thm:closed-split} (with $K=Y\cap Z$, $P=\cl E\cap Y\cap Z$, $Q=\cl F\cap
Y\cap Z$) there are disjoint closed sets $H,L$ with $H\cup L=Y\cap Z$, $\cl E\cap Y\cap Z\subseteq H$ and
$\cl F\cap Y\cap Z\subseteq L$.

\emph{Step 3.} The closed sets $H\cup\cl E$ and $L\cup\cl F$ cover $Z$: a point of $Z\cap Y$ is in
$H\cup L$, and a point of $Z\setminus Y=V$ is in $E\cup F$. They have no common point $x\in Z$: $H\cap
L=\emptyset$; a point of $H\cap\cl F$ lies in $\cl F\cap Y\cap Z\subseteq L$, impossible, and similarly
for $\cl E\cap L$; and a point $x\in\cl E\cap\cl F\cap Z$ either lies in $Y$, and then in $H\cap L$, or
lies in $V=E\cup F$, and then in $E\cap\cl F$ or in $\cl E\cap F$, both empty. Since $Z$ is
connected, it lies in one of the two sets, say $Z\subseteq H\cup\cl E$. Take $b\in F\subseteq Z$; then
$b\notin H\subseteq Y$ and $b\notin\cl E$, a contradiction. The other case is symmetric.
\end{proof}

\formal{\lean{HPC.Plane.isConnected_diff_of_not_weakTriod} (Nadler 11.27), with the two halves
in \lean{HPC.Plane.isConnected_diff_left_of_not_weakTriod}; the continuum $\cl U$ is
\lean{HPC.Plane.isContinuum_closure_diff}.}

\begin{lemma}\label{lem:1128}
Let $X=Y\oplus Z$ not be a weak triod, and let $H$ be a subcontinuum of $\cl U$ with $\cl U\cap
Z\subseteq H$ and $U\setminus H\ne\emptyset$. Then $U\setminus H$ is connected.
\end{lemma}

\begin{proof}
$Z\cup H$ is a continuum: by Lemma~\ref{lem:esum-basic}(c) there is $w\in\cl U\cap Z\subseteq H$. We
claim $X=(Z\cup H)\oplus\cl U$. The union is $X$ because $X=\cl U\cup Z$ and $H\subseteq\cl U\subseteq X$. A
point of $U\setminus H$ is in neither $Z$ nor $H$, so $Z\cup H\ne X$. And $\cl U\neq X$ by
Lemma~\ref{lem:esum-basic}(b). By Theorem~\ref{thm:diff-connected} the set $X\setminus(Z\cup
H)=U\setminus H$ is connected.
\end{proof}

\formal{\lean{HPC.Plane.isConnected_diff_of_superset'} (Nadler 11.28, corrected).}

\begin{remark}\label{rem:1128-false}
Nadler states Lemma~\ref{lem:1128} without the hypothesis $U\setminus H\neq\emptyset$, and then it is
false: for the arc $[0,1]=[0,\frac12]\oplus[\frac12,1]$ (not a weak triod, being irreducible) and
$H=\cl U=[0,\frac12]$, the set $U\setminus H$ is empty, hence not connected. His proof tacitly
assumes the extra hypothesis. Formally: \lean{HPC.Plane.isConnected_diff_of_superset_false}.
\end{remark}

\section{Three lemmas on irreducibility}

\begin{proposition}\label{prop:1131}
Let $X=Y\oplus Z$ not be a weak triod. Let $A,B$ be continua with $A\subseteq\cl U$ and $\cl U\cap
Z\subseteq A$, let $p\in U\setminus A$, let $B=\irr(\set p,A)$, and suppose $U\subseteq A\cup B$. Then
$\cl U=\irr(\set p,\cl U\cap Z)$.
\end{proposition}

\begin{proof}
$\cl U$ is a continuum (Theorem~\ref{thm:diff-connected}) containing $p$ and meeting $\cl U\cap Z$
(Lemma~\ref{lem:esum-basic}(c)). Let $L\subseteq\cl U$ be a subcontinuum containing $p$ and a point
$w\in\cl U\cap Z$; note $w\in A$. We must show $L=\cl U$.

\emph{First, $B\subseteq L$.} Apply Corollary~\ref{cor:bump-open} to $L$ and the open set $O=S\setminus
A$: $p\in L\setminus A$ and $w\in L\cap A$, so there is a subcontinuum $K$ of $L$ with $p\in
K\subseteq\cl{L\setminus A}$ and $K\cap A\ne\emptyset$. Every point $y\in L\setminus A$ lies in $B$: it
is not in $Z$ (otherwise $y\in\cl U\cap Z\subseteq A$), so $y\in U\subseteq A\cup B$, so $y\in B$. As $B$ is
closed, $K\subseteq\cl{L\setminus A}\subseteq B$. Then $K$ is a subcontinuum of $B$ meeting $\set p$ and
$A$, so $K=B$, and $B=K\subseteq L$.

\emph{Second, $L=\cl U$.} Otherwise $U\not\subseteq L$ (else $\cl U\subseteq L$). Take $x\in U\setminus L$;
then $x\in A$, because $x\in A\cup B$ and $B\subseteq L$. Take $q\in Z\setminus Y$
(Lemma~\ref{lem:esum-basic}(a)). Then $X=A\oplus L\oplus Z$ is a weak triod: the union is $X$ (a point
of $X$ is in $Z$ or in $U\subseteq A\cup B\subseteq A\cup L$); $w$ is a common point; $x\in A$ is not in
$L$ nor in $Z$; $p\in L$ is not in $A$ nor in $Z$; and $q\in Z$ is not in $A\cup L\subseteq\cl U$, because
$\cl U$ misses $Z\setminus Y$. This is excluded.
\end{proof}

\formal{\lean{HPC.Plane.nadler_11_31} (Nadler 11.31, without the unused hypothesis $B\subseteq\cl
U$).}

\begin{proposition}\label{prop:1132}
Let $X=Y\oplus Z$ be unicoherent and not a weak triod. Let $A$ be a subcontinuum of $\cl U$ with
$\cl U\cap Z\subseteq A$, and put $W=U\setminus A$. If $\cl W=E\oplus F$, then one of $E,F$ misses $A$
and the other meets $A$.
\end{proposition}

\begin{proof}
$Z\cup A$ is a continuum (it contains a point of $\cl U\cap Z$), and $X=\cl W\cup(Z\cup A)$: a point of
$X$ not in $Z\cup A$ lies in $U\setminus A=W$. Fix $q\in Z\setminus Y$; then $q\notin\cl U\supseteq\cl W$.
Points of $W$ are not in $Z\cup A$. Since $E\neq\cl W$ and $E$ is closed, $W\not\subseteq E$; so there is
$y\in W\setminus E$, and then $y\in F$. Likewise there is $x\in W\setminus F$, and $x\in E$.

\emph{(*) $E\cap F\cap(Z\cup A)=\emptyset$.} A common point $c$ would make $X=E\oplus F\oplus(Z\cup A)$
a weak triod: $x\in E$ is in neither $F$ nor $Z\cup A$; $y\in F$ is in neither $E$ nor $Z\cup A$;
$q\in Z\cup A$ is in neither $E$ nor $F$.

By unicoherence $D=\cl W\cap(Z\cup A)$ is connected. It lies in $E\cup F$, and by (*) it splits into
the disjoint closed sets $D\cap E$ and $D\cap F$; so it lies in one of them, say $D\subseteq E$. Then
$F\cap(Z\cup A)=\emptyset$, since a point of it would lie in $D\subseteq E$, contradicting (*). In
particular $F\cap A=\emptyset$. Finally $D\ne\emptyset$ (as $X=\cl W\cup(Z\cup A)$ is connected and both
sets are closed and nonempty), so $E$ meets $Z\cup A$ at some point $e\in\cl W\subseteq\cl U$; if $e\in Z$
then $e\in\cl U\cap Z\subseteq A$. So $E\cap A\ne\emptyset$.
\end{proof}

\formal{\lean{HPC.Plane.nadler_11_32} (Nadler 11.32, without the hypothesis $A\ne\cl U$, which
follows from the others).}

\begin{lemma}\label{lem:1133a}
Let $X=Y\oplus Z$ be unicoherent, $p\in U$, $q\in V$, and suppose $\cl U=\irr(\set p,\cl U\cap Z)$. If
$K\subseteq X$ is a subcontinuum containing $p$ and $q$, then $U\subseteq K$.
\end{lemma}

\begin{proof}
$K\cup Z$ is a continuum ($q\in K\cap Z$), and $X=(K\cup Z)\cup\cl U$. By unicoherence
$L=(K\cup Z)\cap\cl U$ is connected; it is closed and contains $p$, so it is a subcontinuum of $\cl U$.
It meets $\cl U\cap Z$: otherwise $L\cap Z=\emptyset$, and $K\cup Z\subseteq L\cup Z$ (a point of $K$ not in
$Z$ lies in $U\subseteq\cl U$, hence in $L$), splitting the connected set $K\cup Z$ into the disjoint
closed sets $L\ni p$ and $Z\ni q$. By irreducibility $L=\cl U$. Hence $U\subseteq L\subseteq K\cup Z$, and since
$U\cap Z=\emptyset$, $U\subseteq K$.
\end{proof}

\formal{\lean{HPC.Plane.diff_subset_of_irr}.}

\begin{proposition}\label{prop:1133}
Let $X=Y\oplus Z$ be unicoherent and not a weak triod, $p\in U$, $q\in V$, and suppose
$\cl U=\irr(\set p,\cl U\cap Z)$ and $\cl V=\irr(\set q,\cl V\cap Y)$. Then $X$ is irreducible between
$p$ and $q$.
\end{proposition}

\begin{proof}
Let $K\subseteq X$ be a subcontinuum containing $p$ and $q$, and suppose $K\ne X$. By
Lemma~\ref{lem:1133a} and its mirror image, $U\subseteq K$ and $V\subseteq K$, hence $\cl U\subseteq K$ and
$\cl V\subseteq K$. Note $q\notin\cl U$ (as $\cl U\subseteq Y$).

Apply Corollary~\ref{cor:bump-open} to $K$ and $O=S\setminus\cl U$: there is a subcontinuum $J_0$ of
$K$ with $q\in J_0\subseteq\cl{K\setminus\cl U}$ meeting $\cl U$ at some point $w$. \emph{$J_0$ misses
$U$:} if $x\in J_0\cap U$, the open neighbourhood $S\setminus Z$ of $x$ contains a point $y\in
K\setminus\cl U$, and then $y\in X\setminus Z=U\subseteq\cl U$, which is absurd. Put $J=J_0\cup\cl V$, a
continuum (Theorem~\ref{thm:diff-connected}; both contain $q$); $J$ misses $U$ because $\cl V\subseteq Z$.
$Y\cap Z$ is a continuum by unicoherence.

Take $x\in X\setminus K$. Then $x\in Z$ (else $x\in U\subseteq K$) and $x\in Y$ (else $x\in V\subseteq K$). Also
$w\in X$, $w\notin U$, so $w\in Z$; and $w\in\cl U\subseteq Y$.

Now $X=J\oplus\cl U\oplus(Y\cap Z)$ is a weak triod. The union is $X$: a point of $Z\setminus Y$ is in
$V\subseteq J$, a point of $Z\cap Y$ is in $Y\cap Z$, and a point not in $Z$ is in $U$. The point $w$ is
common to the three sets. $q\in J$ is not in $\cl U$ and not in $Y$. $p\in\cl U$ is not in $J$ (which
misses $U$) and not in $Z$. $x\in Y\cap Z$ is not in $J\cup\cl U\subseteq K$. This is excluded.
\end{proof}

\formal{\lean{HPC.Plane.nadler_11_33} (Nadler 11.33).}

\section{Sorgenfrey's theorem}

The heart of the proof is the following theorem.

\begin{theorem}\label{thm:irr-point}
Let $X=Y\oplus Z$ be unicoherent and not a weak triod. Then there is $p\in U$ with $\cl
U=\irr(\set p,\cl U\cap Z)$.
\end{theorem}

The proof is by contradiction, and occupies the rest of this section. Suppose that \emph{no}
$p\in U$ has $\cl U=\irr(\set p,\cl U\cap Z)$. Write $T=\cl U$, a continuum
(Theorem~\ref{thm:diff-connected}) contained in $Y$, and fix $w\in T\cap Z$ and $q\in Z\setminus Y$
(Lemma~\ref{lem:esum-basic}). Note that $T\subseteq Y$ misses $q$.

\begin{lemma}\label{lem:sorg-proper}
Every $p\in U$ lies in a subcontinuum $L$ of $T$ with $L\ne T$ and $L\cap Z\ne\emptyset$.
\end{lemma}

\begin{proof}
$T$ contains $p$ and meets $T\cap Z$, so, since $T\ne\irr(\set p,T\cap Z)$, some subcontinuum $L\ne T$
of $T$ contains $p$ and meets $T\cap Z$.
\end{proof}

\formal{\lean{HPC.Plane.exists_proper_of_not_irr}.}

Call a set $K$ \emph{good} if $K$ is a subcontinuum of $T$, $K\neq T$, and $T\cap Z\subseteq K$.

\begin{lemma}\label{lem:sorg-step}
\begin{enumerate}[label=(\alph*)]
\item $T\cap Z$ is good.
\item If $K$ is good and $d\in U$, there is a good $K'\supseteq K$ with $d\in K'$.
\end{enumerate}
\end{lemma}

\begin{proof}
(a) $X=T\cup Z$ (Lemma~\ref{lem:esum-basic}), so $T\cap Z$ is a continuum by unicoherence; it misses
$U\ne\emptyset$, while $U\subseteq T$, so $T\cap Z\neq T$.

(b) If $d\in K$ take $K'=K$. Otherwise Theorem~\ref{thm:irr-exists} gives a subcontinuum
$C=\irr(\set d,K)$ of $T$. Put $K'=K\cup C$, a continuum contained in $T$ and containing $T\cap Z$ and
$d$. If $K'=T$, then $U\subseteq K\cup C$, and Proposition~\ref{prop:1131} (with $A=K$, $B=C$, $p=d$)
gives $T=\irr(\set d,T\cap Z)$, contrary to our supposition. So $K'$ is good.
\end{proof}

\formal{\lean{HPC.Plane.SorgGood}, \lean{HPC.Plane.sorgGood_base}, \lean{HPC.Plane.sorgGood_step}.}

The compact metric space $S$ is separable, and so is its subset $U$: choose $d_0,d_1,\dots\in U$ with
$U\subseteq\cl{\set{d_k:k\in\N}}$. Using Lemma~\ref{lem:sorg-step} define good sets
\[
K_0=T\cap Z\subseteq K_1\subseteq K_2\subseteq\cdots,\qquad d_n\in K_{n+1}.
\]
(Nadler absorbs at each stage the first $d_k$ not yet absorbed; absorbing $d_n$ at stage $n$ is
simpler and suffices.) Put
\[
U_n=U\setminus K_n .
\]

\begin{lemma}\label{lem:sorg-facts}
For every $n$:
\begin{enumerate}[label=(\alph*)]
\item $U_n\ne\emptyset$, $\cl{U_n}$ is a continuum contained in $T$, and $U_{n+1}\subseteq U_n$;
\item $X=Z\cup K_n\cup U_n$ and $\cl{U_n}\subseteq U_n\cup K_n$;
\item $\cl{U_n}\cap K_n$ is a continuum;
\item $\cl{U_n}$ is decomposable.
\end{enumerate}
\end{lemma}

\begin{proof}
(a) If $U_n=\emptyset$ then $U\subseteq K_n$ and $T=\cl U\subseteq K_n$, contradicting $K_n\ne T$. By
Lemma~\ref{lem:1128} (with $H=K_n$) $U_n$ is connected, so $\cl{U_n}$ is a continuum; it lies in
$T$. The $K_n$ increase, so the $U_n$ decrease.

(b) A point of $X$ not in $Z$ is in $U$, hence in $K_n$ or $U_n$. A point $x\in\cl{U_n}$ not in $K_n$ is
not in $Z$ (else $x\in T\cap Z\subseteq K_n$), so $x\in U\setminus K_n=U_n$.

(c) $Z\cup K_n$ is a continuum ($w$ is in both), and $X=\cl{U_n}\cup(Z\cup K_n)$ by (b). By
unicoherence $\cl{U_n}\cap(Z\cup K_n)$ is connected, and it equals $\cl{U_n}\cap K_n$ because
$\cl{U_n}\cap Z\subseteq T\cap Z\subseteq K_n$. It is closed, so it is a continuum.

(d) Suppose $I=\cl{U_m}$ is indecomposable. $P=I\cap K_m$ is a continuum by (c); take $a\in P$. By
\vtwo{Theorem~1.4} there is $r\in I$ lying in no proper subcontinuum of $I$ containing $a$. Fix $u\in
U_m$; then $u\in I\setminus P$, so $P\ne I$.

\emph{If $C\subseteq I$ is a subcontinuum containing $r$ and meeting $P$, then $C=I$.} Indeed $C\cup P$ is a
continuum containing $a$ and $r$, so $C\cup P=I$; by indecomposability $C=I$ or $P=I$, and $P\neq I$.

\emph{$r\notin K_m$}: otherwise $C=\set r$ meets $P$, so $\set r=I$, and $u=r\in K_m$, which is false.
By (b), $r\in U_m\subseteq U$. By Lemma~\ref{lem:sorg-proper} there is a subcontinuum $L\ne T$ of $T$
containing $r$ and a point $z\in Z$. Apply Corollary~\ref{cor:bump-open} to $L$ and $O=S\setminus(Z\cup
K_m)$ (note $r\in O$, $z\notin O$): there is a subcontinuum $C$ of $L$ with $r\in C\subseteq\cl{L\cap O}$
meeting $Z\cup K_m$ at some point $x$. Now $L\cap O\subseteq U_m$ (a point of $L\subseteq X$ outside $Z$ is
in $U$), so $C\subseteq I$; and $x\in P$ (if $x\in Z$ then $x\in T\cap Z\subseteq K_m$). So $C=I$, and
$I\subseteq L$.

$K_m\not\subseteq Z\cup L$: otherwise every point of $U$ lies in $L$ (a point of $U\cap K_m$ is in $L$ as
it is not in $Z$, and $U_m\subseteq I\subseteq L$), so $T=\cl U\subseteq L$, contradicting $L\ne T$. Take $k\in
K_m\setminus(Z\cup L)$. Then $X=Z\oplus L\oplus K_m$ is a weak triod: the union is $X$ by (b), since
$U_m\subseteq I\subseteq L$; the point $z$ lies in $Z\cap L$ and in $T\cap Z\subseteq K_m$; $q\in Z$ is not in
$L\cup K_m\subseteq T$; $u\in L$ is not in $Z\cup K_m$; $k\in K_m$ is not in $Z\cup L$. This is excluded.
\end{proof}

\begin{lemma}\label{lem:sorg-star}
For every $\ell$ there is $n>\ell$ with $K_\ell\cap\cl{U_n}=\emptyset$.
\end{lemma}

\begin{proof}
Suppose that $K_\ell\cap\cl{U_n}\neq\emptyset$ for all $n>\ell$. By Lemma~\ref{lem:sorg-facts}(d)
write $\cl{U_\ell}=E\oplus F$. By Proposition~\ref{prop:1132} (with $A=K_\ell$, $W=U_\ell$) one of $E,F$
meets $K_\ell$ and the other misses it; by symmetry say $E\cap K_\ell\ne\emptyset$ and $F\cap
K_\ell=\emptyset$. Then $F\subseteq U_\ell$ by Lemma~\ref{lem:sorg-facts}(b), and $U_\ell\subseteq E\cup F$.

Since $E\ne\cl{U_\ell}$, there is $y\in\cl{U_\ell}\setminus E$; then $y\in F\subseteq U_\ell$. The open set
$S\setminus(E\cup K_\ell)$ contains $y\in U\subseteq\cl{\set{d_k}}$, so it contains some $d_k$. Then
$d_k\in U\setminus K_\ell=U_\ell$ and $d_k\notin E$, so $d_k\in F$. Let $j=\max(k+1,\ell+1)>\ell$; then
$d_k\in K_{k+1}\subseteq K_j$ and $K_\ell\subseteq K_j$. Fix $a\in K_\ell$.

Apply Corollary~\ref{cor:bump-open} to $K_j$ and $O=S\setminus F$ ($a\in K_j\setminus F$, $d_k\in K_j\cap
F$): there is a subcontinuum $J_0$ of $K_j$ with $a\in J_0\subseteq\cl{K_j\setminus F}$ meeting $F$ at some
point $f$. Put $J=J_0\cup K_\ell$, a continuum contained in $K_j$.

\emph{Every point of $J\cap F$ lies in $E$.} Points of $K_\ell$ are not in $F$. Let $x\in J_0\cap F$
with $x\notin E$. Then $x\in U_\ell$, so $x\notin Z\cup K_\ell$, and the open set $S\setminus(Z\cup
K_\ell\cup E)$ contains $x\in\cl{K_j\setminus F}$, hence a point $y'\in K_j\setminus F$. This $y'$ is in
$U$ (not in $Z$), not in $K_\ell$, so in $U_\ell\subseteq E\cup F$; but it is in neither $E$ nor $F$.

The sets $Z\cup J$, $E\cup J$ and $F\cup J$ are continua (they meet at $w\in K_\ell$, at a point of
$E\cap K_\ell$, and at $f$, respectively).

\emph{Some point $z\in E$ lies outside $Z$, $J$ and $F$.} Otherwise every point of $E\setminus F$ lies in
$K_j$ (a point of $E\cap Z$ lies in $T\cap Z\subseteq K_j$, and $J\subseteq K_j$). Then $U_j\subseteq F$: a point
of $U_j\subseteq U_\ell\subseteq E\cup F$ in $E\setminus F$ would be in $K_j$, but $U_j$ misses $K_j$. So
$\cl{U_j}\subseteq F$, and by our supposition ($j>\ell$) $K_\ell$ meets $\cl{U_j}\subseteq F$, contradicting
$F\cap K_\ell=\emptyset$.

Now $X=(Z\cup J)\oplus(E\cup J)\oplus(F\cup J)$ is a weak triod. The union is $X$ by
Lemma~\ref{lem:sorg-facts}(b), as $K_\ell\subseteq J$ and $U_\ell\subseteq E\cup F$. The point $a$ is common
to the three sets. A point $v\in V$ lies in $Z$ but not in $T\supseteq E\cup F\cup J$. The point $z$ lies in
$E$ but not in $Z\cup J$ nor $F\cup J$. The point $y$ lies in $F$, not in $Z$ (as $y\in U$), not in $E$,
and not in $J$ (a point of $J\cap F$ lies in $E$). This is excluded.
\end{proof}

\begin{proof}[Proof of Theorem~\ref{thm:irr-point}]
The continua $\cl{U_n}$ decrease, so by Cantor's theorem there is $t\in\bigcap_n\cl{U_n}$. By
Lemma~\ref{lem:sorg-star} $t$ lies in no $K_\ell$. By Lemma~\ref{lem:sorg-facts}(b) (with $n=0$),
$t\in U$. By Lemma~\ref{lem:sorg-proper} there is a subcontinuum $Q\ne T$ of $T$ containing $t$ and a
point $y\in Z$.

\emph{Some $K_\ell$ has a point $x\notin Z\cup Q$.} Otherwise every $d_k$ (which lies in $K_{k+1}$ and
not in $Z$) lies in $Q$, so $U\subseteq\cl{\set{d_k}}\subseteq Q$ and $T\subseteq Q$, contradicting $Q\neq T$.

By Lemma~\ref{lem:sorg-star} take $n>\ell$ with $K_\ell\cap\cl{U_n}=\emptyset$; so $x\in K_n$ and
$x\notin\cl{U_n}$. The set $Q\cup\cl{U_n}$ is a continuum (both contain $t$). Now
$X=Z\oplus K_n\oplus(Q\cup\cl{U_n})$ is a weak triod: the union is $X$ by
Lemma~\ref{lem:sorg-facts}(b); $y\in Z\cap Q$, and $y\in T\cap Z\subseteq K_n$, is a common point; $q\in Z$
is not in $K_n\cup Q\cup\cl{U_n}\subseteq T$; $x\in K_n$ is not in $Z$ nor in $Q\cup\cl{U_n}$; and $t\in Q$
is not in $Z$ nor in $K_n$. This contradiction proves the theorem.
\end{proof}

\formal{\lean{HPC.Plane.exists_irrFromTo_point} (Nadler 11.34(a)); the body of the proof by
contradiction is \lean{HPC.Plane.sorgenfrey_contra}.}

\begin{theorem}[Sorgenfrey]\label{thm:sorgenfrey}
A nondegenerate unicoherent continuum that is not a triod is irreducible between two (distinct)
points.
\end{theorem}

\begin{proof}
Let $X$ be such a continuum. If $X$ is indecomposable, apply Theorem~\ref{thm:indec-irr}. Otherwise
$X=Y\oplus Z$ for some subcontinua $Y,Z$. By Proposition~\ref{prop:uni-triod} $X$ is not a weak triod.
Theorem~\ref{thm:irr-point}, applied to $X=Y\oplus Z$ and to $X=Z\oplus Y$, gives $p\in U$ with $\cl
U=\irr(\set p,\cl U\cap Z)$ and $q\in V$ with $\cl V=\irr(\set q,\cl V\cap Y)$. By
Proposition~\ref{prop:1133} $X$ is irreducible between $p$ and $q$; and $p\ne q$ because $p\notin Z$
and $q\in V\subseteq Z$.
\end{proof}

\formal{\lean{HPC.Plane.sorgenfrey} (Nadler 11.34).}

\begin{corollary}\label{cor:irr-iff}
A nondegenerate unicoherent continuum is irreducible if and only if it is not a triod.
\end{corollary}

\begin{proof}
A triod is a weak triod (Proposition~\ref{prop:triod-weak}), and a weak triod is not irreducible
(Lemma~\ref{lem:weak-not-irr}). The converse is Theorem~\ref{thm:sorgenfrey}.
\end{proof}

\formal{\lean{HPC.Plane.irreducible_iff_not_isTriod} (Nadler 11.35).}

\begin{corollary}\label{cor:atriodic-irr}
A nondegenerate, atriodic, hereditarily unicoherent continuum is irreducible between two of its
points.
\end{corollary}

\formal{\lean{HPC.Plane.exists_irreducibleBetween_of_atriodic}.}

\begin{remark}
Corollary~\ref{cor:atriodic-irr} is the form used for snake-like continua (which are atriodic and
hereditarily unicoherent by Chapter~\ref{ch:snake}); with it, the formal proof derives Fugate's
characterization of terminal subcontinua (\lean{HPC.Plane.isTerminalSub_iff_irreducible}).
\end{remark}

\section*{Exercises}

\begin{exercise}
Show that the union of three arcs with a common end point and otherwise disjoint (a simple triod)
is a triod, is unicoherent, and is not irreducible.
\end{exercise}

\begin{exercise}
Show that the $\sin(1/x)$-continuum is unicoherent and atriodic, and find two points between which
it is irreducible.
\end{exercise}

\begin{exercise}
Show that the hypothesis of unicoherence cannot be dropped from Theorem~\ref{thm:sorgenfrey}: a
circle is not a triod and not irreducible.
\end{exercise}

%======================================================================
\chapter{Snake-like and circle-like continua}\label{ch:snake}
%======================================================================

Recall from Volume~1 that a chain $D=(D_0,\dots,D_{n-1})$ is a sequence of open sets with $D_i\cap
D_j\ne\emptyset$ if and only if $\abs{i-j}\le1$ (\vone{Definition~2.1}); it is an
\emph{$\varepsilon$-chain} if every link is bounded with diameter $<\varepsilon$
(\vone{Definition~2.3}); and a continuum $M\subseteq X$ is \emph{snake-like} if for every $\varepsilon>0$
some $\varepsilon$-chain covers it (\vone{Definition~2.13}). We use repeatedly:

\begin{quote}
\vone{Lemma~2.8} (separation by a link). \emph{If $C\subseteq\bigcup E$ is preconnected and meets the
links $E_a$ and $E_c$ of a chain $E$, and $a<b<c$, then $C$ meets $E_b$.}
\end{quote}

\section{Subcontinua}

\begin{theorem}\label{thm:snake-sub}
A subcontinuum of a snake-like continuum is snake-like.
\end{theorem}

\begin{proof}
Let $K\subseteq M$ with $M$ snake-like, and $\varepsilon>0$. Take an $\varepsilon$-chain $D$ of length $n$
covering $M$. Let $i$ and $j$ be the least and the largest index of a link meeting $K$ (there is one,
as $K\ne\emptyset$). By \vone{Lemma~2.8} every link $D_l$ with $i\le l\le j$ meets $K$, and the others do
not. The links $D_i,D_{i+1},\dots,D_j$ form a chain (the chain condition only involves differences
of indices), an $\varepsilon$-chain, which covers $K$.
\end{proof}

\formal{\lean{HPC.Plane.isSnakeLike_of_subcontinuum}; the interval of links meeting $K$ is
\lean{HPC.Plane.meets_between}.}

\section{Unicoherence and triods}

\begin{lemma}[Intermediate values along a chain]\label{lem:chain-ivt}
Let $D$ be an $\varepsilon$-chain and $f\colon X\to\R$ a function with $\abs{f(x)-f(y)}\le d(x,y)$.
Let $i\le j$ be indices and $c\in\R$, and suppose the link $D_i$ contains a point with $f\le c$ and
$D_j$ a point with $f\ge c$. Let $Y$ be a set meeting every link $D_l$, $i\le l\le j$. Then there are
$l\in[i,j]$ and $y\in Y\cap D_l$ with $\abs{f(y)-c}<\varepsilon$. The same holds if $i>j$, with $[j,i]$
in place of $[i,j]$.
\end{lemma}

\begin{proof}
The case $i>j$ follows by reversing the chain. Let $l$ be the largest index in $[i,j]$ such that $D_l$
contains a point $x$ with $f(x)\le c$. If $l=j$, then $D_l$ also contains a point $x'$ with $f(x')\ge c$.
If $l<j$, every point of $D_{l+1}$ has $f>c$, and $D_l\cap D_{l+1}$ contains such a point $x'$. In both
cases $D_l$ contains points $x,x'$ with $f(x)\le c\le f(x')$. Take $y\in Y\cap D_l$. If $f(y)\ge c$ then
$0\le f(y)-c\le f(y)-f(x)\le d(y,x)<\varepsilon$; if $f(y)<c$ then $0<c-f(y)\le f(x')-f(y)<\varepsilon$.
\end{proof}

\formal{\lean{HPC.Plane.chain_ivt}.}

\begin{corollary}\label{cor:two-sets}
Under the hypotheses of Lemma~\ref{lem:chain-ivt}, if $Y$ and $Z$ both meet every link between $D_i$
and $D_j$, there are $y\in Y$ and $z\in Z$ with $\abs{f(y)-c}<\varepsilon$ and $d(y,z)<\varepsilon$.
\end{corollary}

\begin{proof}
Take $l$ and $y$ from the lemma and $z\in Z\cap D_l$; then $d(y,z)\le\diam D_l<\varepsilon$.
\end{proof}

\formal{\lean{HPC.Plane.chain_two_sets}.}

In the next two proofs $S$ is a compact metric space, and $\dist(x,P)$ is the distance from $x$ to a
nonempty set $P$, a function with $\abs{\dist(x,P)-\dist(y,P)}\le d(x,y)$. We also use: \emph{two
disjoint compact sets $A,B$ have $d(a,b)\ge\eta$ for all $a\in A$, $b\in B$, for some $\eta>0$}
(\vone{Lemma~1.1}, or trivially if one of them is empty).

\begin{theorem}\label{thm:snake-unicoherent}
A snake-like continuum in a compact metric space is unicoherent, and hence hereditarily
unicoherent.
\end{theorem}

\begin{proof}
Let $M$ be snake-like and $M=H\cup K$ with $H,K$ subcontinua. Since $M$ is connected and $H,K$ are
closed and nonempty, $H\cap K\neq\emptyset$. Suppose $H\cap K$ is not connected: $H\cap K=P\cup Q$ with
$P,Q$ closed, disjoint and nonempty. Let $\delta>0$ be the distance from $Q$ to $P$, and
$f(x)=\dist(x,P)$; so $f=0$ on $P$ and $f\ge\delta$ on $Q$. Consider the band
\[
B=\set{x:\delta/4\le f(x)\le3\delta/4}.
\]
The compact sets $H\cap B$ and $K\cap B$ are disjoint, because a common point would lie in $H\cap
K=P\cup Q$, where $f=0$ or $f\ge\delta$. Let $\eta>0$ be a lower bound for their mutual distances, and
$\varepsilon=\min(\eta,\delta/8)$.

Take an $\varepsilon$-chain $D$ covering $M$, and links $D_i\ni p$ and $D_j\ni q$ with $p\in P$, $q\in
Q$. The connected sets $H$ and $K$ contain $p$ and $q$, so by \vone{Lemma~2.8} each meets every link
between $D_i$ and $D_j$. Corollary~\ref{cor:two-sets} with $c=\delta/2$ (note $f(p)=0\le c\le f(q)$)
gives $y\in H$ and $z\in K$ with $\abs{f(y)-\delta/2}<\delta/8$ and $d(y,z)<\varepsilon$; then also
$\abs{f(z)-\delta/2}<\delta/4$. So $y\in H\cap B$, $z\in K\cap B$, and $d(y,z)\ge\eta\ge\varepsilon$, a
contradiction.

Every subcontinuum of $M$ is snake-like (Theorem~\ref{thm:snake-sub}), hence unicoherent.
\end{proof}

\formal{\lean{HPC.Plane.isUnicoherent_of_isSnakeLike} and
\lean{HPC.Plane.isHereditarilyUnicoherent_of_isSnakeLike} (Nadler 12.2).}

\begin{theorem}\label{thm:snake-atriodic}
A snake-like continuum in a compact metric space is not a triod. Hence it is atriodic.
\end{theorem}

\begin{proof}
Let $T$ be snake-like with core $N$ and $T\setminus N=A\cup B\cup C$ as in Definition~\ref{def:triod}.
By Lemma~\ref{lem:piece} the sets $N\cup A$, $N\cup B$, $N\cup C$ are continua. Let
$f(x)=\dist(x,N)$, choose $a\in A$, $b\in B$, $c\in C$ (they are not in $N$, so $f>0$ there) and let
$\rho=\min(f(a),f(b),f(c))>0$. The band $B'=\set{x:\rho/4\le f(x)\le3\rho/4}$ misses $N$.

The sets $A\cap B'$, $B\cap B'$, $C\cap B'$ are closed: a point of $\cl{A\cap B'}$ lies in $T$ and in
$B'$, hence not in $N$, and not in $B\cup C$ (which is separated from $A$), so in $A\cap B'$. They are
pairwise disjoint, so there is $\eta>0$ bounding their mutual distances from below. Let
$\varepsilon=\min(\eta,\rho/8)$, and take an $\varepsilon$-chain covering $T$ with links $D_{i_a}\ni a$,
$D_{i_b}\ni b$, $D_{i_c}\ni c$ and $D_{i_n}\ni n$ for some $n\in N$.

Two of the three indices $i_a,i_b,i_c$ lie on the same side of $i_n$; say (renaming) $i_a$ and $i_b$,
with $i_a$ not farther from $i_n$ than $i_b$. Then every index between $i_n$ and $i_a$ lies between
$i_n$ and $i_b$. The connected set $N\cup A$ contains $n$ and $a$, and $N\cup B$ contains $n$ and $b$,
so both meet every link between $D_{i_n}$ and $D_{i_a}$ (\vone{Lemma~2.8}). Apply
Corollary~\ref{cor:two-sets} with the level $\rho/2$ in place of $c$ ($f(n)=0$, $f(a)\ge\rho$): there are $y\in N\cup A$ and
$z\in N\cup B$ with $\abs{f(y)-\rho/2}<\rho/8$ and $d(y,z)<\varepsilon$, so that $y,z\in B'$. Then
$y,z\notin N$, so $y\in A\cap B'$ and $z\in B\cap B'$, and $d(y,z)\ge\eta\ge\varepsilon$, a contradiction.

Every subcontinuum of a snake-like continuum is snake-like, hence not a triod.
\end{proof}

\formal{\lean{HPC.Plane.isAtriodic_of_isSnakeLike}, through
\lean{HPC.Plane.not_isTriod_of_isSnakeLike} (Nadler 12.4).}

\begin{remark}
Bing (1951, Theorem~11) proved the converse for hereditarily decomposable continua: an atriodic,
hereditarily unicoherent, hereditarily decomposable continuum is snake-like. This is not needed
for the classification. By Corollary~\ref{cor:atriodic-irr}, a nondegenerate snake-like continuum
is irreducible between two of its points.
\end{remark}

\section{\texorpdfstring{$\varepsilon$}{Epsilon}-maps}

\begin{definition}\label{def:epsmap}
Let $K\subseteq X$ and $\varepsilon>0$. An \emph{$\varepsilon$-map} on $K$ is a continuous function
$g\colon K\to\R$ such that $g(x)=g(y)$ implies $d(x,y)<\varepsilon$.
\end{definition}

\begin{proposition}\label{prop:chain-epsmap}
If $K$ is covered by an $\varepsilon$-chain $D$, then $K$ has an $\varepsilon$-map.
\end{proposition}

\begin{proof}
For each $i$ let $w_i(x)=\dist(x,X\setminus D_i)$ (or $w_i=1$ if $D_i=X$): a continuous function, $\ge0$,
and $>0$ exactly on the open set $D_i$. Put $\Sigma(x)=\sum_iw_i(x)$ and $\Lambda(x)=\sum_i i\,w_i(x)$; on
$K$, $\Sigma>0$. Define $g=\Lambda/\Sigma$ on $K$, a continuous function.

For $x\in K$ let $p$ be the least index with $x\in D_p$. Any two links containing $x$ meet, so their
indices differ by at most one; hence $x$ lies only in $D_p$ and possibly $D_{p+1}$. Then
\[
\Lambda(x)-p\,\Sigma(x)=\sum_i(i-p)w_i(x)\ge0,\qquad(p+1)\Sigma(x)-\Lambda(x)=\sum_i(p+1-i)w_i(x)\ge
w_p(x)>0,
\]
so $p\le g(x)<p+1$, that is, $\fl{g(x)}=p$. If $g(x)=g(y)$, then $x$ and $y$ have the same least link
$D_p$, and $d(x,y)\le\diam D_p<\varepsilon$.
\end{proof}

\formal{\lean{HPC.BJ59.exists_epsMap_of_chain} and \lean{HPC.BJ59.exists_epsMap_of_isSnakeLike}.}

\begin{lemma}\label{lem:uniform-fibres}
Let $M$ be compact and $f\colon M\to\R$ continuous with $f(x)=f(y)\Rightarrow d(x,y)<\varepsilon$ for
$x,y\in M$. Then there is $\delta>0$ such that $x,y\in M$ and $\abs{f(x)-f(y)}<\delta$ imply
$d(x,y)<\varepsilon$.
\end{lemma}

\begin{proof}
The set $\set{(x,y)\in M\times M:d(x,y)\ge\varepsilon}$ is compact, and the continuous function
$\abs{f(x)-f(y)}$ is positive on it; let $\delta$ be its minimum (or $\delta=1$ if the set is empty).
\end{proof}

\formal{\lean{HPC.Plane.exists_delta_of_fibres_on}.}

\begin{theorem}[The $\varepsilon$-map criterion]\label{thm:epsmap-criterion}
Let $S$ be a compact metric space and $M\subseteq S$ a continuum. Suppose that for every
$\varepsilon>0$ there is a continuous $f\colon S\to\R$ such that $f(x)=f(y)$ implies $d(x,y)<\varepsilon$
for $x,y\in M$. Then $M$ is snake-like.
\end{theorem}

\begin{proof}
Let $\varepsilon>0$, and take $f$ for $\varepsilon/4$. Lemma~\ref{lem:uniform-fibres} gives $\delta>0$
with $\abs{f(x)-f(y)}<\delta\Rightarrow d(x,y)<\varepsilon/4$ on $M$; uniform continuity of $f$ on $S$
gives $\eta>0$ with $d(x,y)<\eta\Rightarrow\abs{f(x)-f(y)}<\delta/6$. Put $\eta'=\min(\eta,\varepsilon/8)$
and let $N$ be the open $\eta'$-neighbourhood of $M$.

$f(M)$ is compact and connected, so $f(M)=[a,b]$. Let $h=\delta/6$ and $n=\fl{(b-a)/h}+1$, so that
$(n-1)h\le b-a$. For $i\in[n]$ put $I_i=(a+(i-1)h,\ a+(i+1)h)$ and
\[
D_i=f^{-1}(I_i)\cap N .
\]
These are open sets. \emph{They form a chain:} if $\abs{i-j}\ge2$ then $I_i\cap I_j=\emptyset$; $D_i$
contains a point of $M$ with $f=a+ih\in[a,b]$; and $D_i\cap D_{i+1}$ contains a point of $M$ with
$f=a+(i+\frac12)h$, which lies in $[a,b]$ since $(i+1)h\le(n-1)h\le b-a$. \emph{They cover $M$:} for
$x\in M$ let $i=\fl{(f(x)-a)/h}$; then $0\le i\le n-1$ and $f(x)\in I_i$. \emph{Their diameters are at
most $\varepsilon/2$:} if $x,y\in D_i$, choose $x',y'\in M$ with $d(x,x'),d(y,y')<\eta'$. Then
$\abs{f(x)-f(y)}<2h=\delta/3$ and $\abs{f(x)-f(x')},\abs{f(y)-f(y')}<\delta/6$, so
$\abs{f(x')-f(y')}<\delta$, hence $d(x',y')<\varepsilon/4$ and $d(x,y)<\varepsilon/4+2\eta'\le\varepsilon/2$.
So $D$ is an $\varepsilon$-chain covering $M$.
\end{proof}

\formal{\lean{HPC.Plane.isSnakeLike_of_epsMap}; for the whole space,
\lean{HPC.Plane.isSnakeLike_univ_of_epsMap}.}

\begin{theorem}[Snake-likeness is topological]\label{thm:snake-homeo}
Let $A\subseteq X$ be snake-like, $Y$ a compact metric space, $B\subseteq Y$, and $e\colon A\to B$ a
homeomorphism. Then $B$ is snake-like.
\end{theorem}

\begin{proof}
$B=e(A)$ is a continuum. Let $\varepsilon>0$. Since $A$ is compact, $e$ is uniformly continuous: there is
$\delta>0$ with $d(a,a')<\delta\Rightarrow d(e(a),e(a'))<\varepsilon$. Proposition~\ref{prop:chain-epsmap}
gives a $\delta$-map $g$ on $A$. The function $g\circ e^{-1}\colon B\to\R$ is continuous on the closed
set $B$; by the Tietze extension theorem it extends to a continuous $F\colon Y\to\R$. If $x,y\in B$ and
$F(x)=F(y)$, then $g(e^{-1}x)=g(e^{-1}y)$, so $d(e^{-1}x,e^{-1}y)<\delta$ and $d(x,y)<\varepsilon$. By
Theorem~\ref{thm:epsmap-criterion}, $B$ is snake-like.
\end{proof}

\formal{\lean{HPC.BJ59.isSnakeLike_of_homeomorph}.}

\begin{remark}
Bing (1951, Theorem~4) proved that every snake-like continuum is homeomorphic to a plane continuum
(\lean{HPC.Plane.theorem4}). The proof combines minimal covering chains with the matching covers of
\vtwo{Chapter~4} and tame chains in the plane; it is not reproduced here.
\end{remark}

\section{Circular chains and circle-like continua}

For $i,j\in[n]$ let the \emph{cyclic distance} be $\min(\abs{i-j},\,n-\abs{i-j})$.

\begin{definition}\label{def:circchain}
A \emph{circular chain} of length $n$ is a sequence $O_0,\dots,O_{n-1}$ of open sets with $O_i\cap
O_j\ne\emptyset$ if and only if the cyclic distance of $i$ and $j$ is at most $1$. It is an
\emph{$\varepsilon$-circular chain} if its links are bounded with diameter $<\varepsilon$. A continuum
is \emph{circle-like} if for every $\varepsilon>0$ some $\varepsilon$-circular chain covers it.
\end{definition}

\formal{\lean{HPC.Plane.CircChain}, \lean{HPC.Plane.IsCircEpsChain}, \lean{HPC.Plane.IsCircleLike}
(in \lean{HPC/Plane/Ingram68.lean}).}

For $n\le3$ every two links of a circular chain meet.

\begin{proposition}\label{prop:circle-circlelike}
The circle $\mathbb T=\set{z\in\mathbb C:\abs z=1}$ is circle-like.
\end{proposition}

\begin{proof}
Let $\varepsilon>0$, choose $n\ge4$ with $3\pi/n<\varepsilon$, put $p(t)=e^{2\pi it/n}$, and
$O_k=p\bigl((k-\frac14,\,k+\frac54)\bigr)$ for $k\in[n]$. These are open arcs (the map $p$ is open). The
parameter intervals have length $\frac32$; two of them meet if and only if their left end points
differ by less than $\frac32$ modulo $n$, that is, if and only if the indices have cyclic distance at
most $1$ (for indices at cyclic distance $\ge2$ the intervals are at least $\frac12$ apart, also
modulo $n$, as $n\ge4$). Every real $t$ lies in $(k-\frac14,k+\frac54)$ for $k=\fl t$, so the arcs
cover $\mathbb T$. Finally $\abs{e^{ia}-e^{ib}}\le\abs{a-b}$, so $\diam O_k\le\frac{2\pi}n\cdot\frac32=
3\pi/n<\varepsilon$.
\end{proof}

\formal{\lean{HPC.BJ59.isCircleLike_circle}, with the arcs \lean{HPC.BJ59.circArc}.}

\begin{proposition}\label{prop:circle-homeo}
Let $X$ be a compact metric space homeomorphic to a circle-like compact metric space $Y$. Then $X$ is
circle-like.
\end{proposition}

\begin{proof}
Let $e\colon X\to Y$ be a homeomorphism; $e^{-1}$ is uniformly continuous. Given $\varepsilon$, take
$\delta$ with $d(y,y')<\delta\Rightarrow d(e^{-1}y,e^{-1}y')<\varepsilon/2$, and a $\delta$-circular
chain $O$ covering $Y$. The sets $e^{-1}(O_i)$ are open, cover $X$, meet exactly when the $O_i$ do
($e$ is a bijection), and have diameter $\le\varepsilon/2$.
\end{proof}

\formal{\lean{HPC.BJ59.isCircleLike_univ_of_homeomorph}.}

\begin{theorem}\label{thm:circle-sub}
A proper subcontinuum of a circle-like continuum is snake-like.
\end{theorem}

\begin{proof}
Let $Z\subsetneq C$ with $C$ circle-like, and $\varepsilon>0$. Take $x\in C\setminus Z$; then
$r=\dist(x,Z)>0$. Take a $\min(\varepsilon,r)$-circular chain $O_0,\dots,O_{n-1}$ covering $C$, and a link
$O_{i_0}\ni x$. This link misses $Z$, since its diameter is less than $r$. List the remaining links in
cyclic order starting after $i_0$: $E_j=O_{\sigma(j)}$ with $\sigma(j)=i_0+1+j$ modulo $n$, for
$j\in[n-1]$. Consecutive $E_j$ are cyclically adjacent, so they meet. If $\abs{j-j'}\ge2$, the cyclic
distance of $\sigma(j)$ and $\sigma(j')$ is $\min(\abs{j-j'},n-\abs{j-j'})\ge2$, because
$\abs{j-j'}\le n-2$; so $E_j\cap E_{j'}=\emptyset$. Thus $E$ is an $\varepsilon$-chain, and it covers $Z$,
as every point of $Z$ lies in a link other than $O_{i_0}$.
\end{proof}

\formal{\lean{HPC.Plane.isSnakeLike_of_subset_circleLike} (Nadler 12.51).}

\section*{Exercises}

\begin{exercise}
Show that the letter T is not snake-like, directly from the definition.
\end{exercise}

\begin{exercise}
Show that the circle is not snake-like. (Hint: Theorem~\ref{thm:snake-unicoherent}.)
\end{exercise}

\begin{exercise}
Show that a continuum $M$ in a compact space $S$ is snake-like if and only if for every
$\varepsilon>0$ it has an $\varepsilon$-map into $[0,1]$.
\end{exercise}

\begin{exercise}
Show that the $\sin(1/x)$-continuum is snake-like, by exhibiting $\varepsilon$-maps and using
Theorem~\ref{thm:epsmap-criterion}.
\end{exercise}

%======================================================================
\chapter{Upper semicontinuous decompositions}\label{ch:decomp}
%======================================================================

Throughout this chapter $X$ is a compact metric space. A \emph{decomposition} of $X$ is a partition
of $X$ into nonempty sets, its \emph{elements}; equivalently, an equivalence relation $r$ on $X$. We
write $\elt x$ for the element containing $x$, so $y\in\elt x$ means $r(y,x)$, and $\elt y=\elt x$ for
every $y\in\elt x$.

\section{Upper and lower semicontinuity}

\begin{definition}\label{def:usc}
A decomposition is \emph{upper semicontinuous} (u.s.c.) if every element is closed and, whenever
an element $\elt x$ lies in an open set $U$, there is an open set $V\supseteq\elt x$ such that every
element meeting $V$ lies in $U$. It is \emph{lower semicontinuous at the element $\elt x$} if for
every $y\in\elt x$ and every open $V\ni y$ there is an open $W\supseteq\elt x$ such that every element
meeting $W$ meets $V$. It is \emph{continuous} if it is u.s.c.\ and lower semicontinuous at every
element.
\end{definition}

\formal{\lean{HPC.Jones.IsUSCDecomp}, \lean{HPC.Jones.LSCAt}, \lean{HPC.Jones.IsContinuousDecomp},
with the element $\elt x$ written \lean{HPC.Jones.cls} (in \lean{HPC/Jones/Decomp.lean}).}

Informally: u.s.c.\ means elements near $\elt x$ lie in any given neighbourhood of $\elt x$ (they
cannot ``suddenly grow''); lower semicontinuity at $\elt x$ means elements near $\elt x$ come near
every point of $\elt x$ (they cannot ``suddenly shrink'').

\begin{example}
The decomposition of $[0,1]^2$ into vertical segments $\set t\times[0,1]$ is continuous. The
decomposition of $[0,1]^2$ into the segment $\set0\times[0,1]$ and the single points not on it is
u.s.c.\ but not lower semicontinuous at the segment.
\end{example}

\begin{lemma}\label{lem:usc-open}
If the decomposition is u.s.c.\ and $U$ is open, then $U^*=\set{x:\elt x\subseteq U}$ is open. It is a
union of elements.
\end{lemma}

\begin{proof}
If $\elt x\subseteq U$, take $V$ as in the definition; every $z\in V$ has $\elt z\subseteq U$, so $V\subseteq
U^*$, and $x\in V$. The second statement holds because $\elt z=\elt x$ for $z\in\elt x$.
\end{proof}

\formal{\lean{HPC.Jones.isOpen_setOf_cls_subset}.}

\begin{proposition}\label{prop:closed-usc}
If the relation $r$ is a closed subset of $X\times X$, the decomposition is u.s.c.
\end{proposition}

\begin{proof}
$\elt x$ is the preimage of the closed set $r$ under $y\mapsto(y,x)$, so it is closed. Let $\elt x\subseteq U$
with $U$ open, and let $V=U^*$. Then $V\supseteq\elt x$, every element meeting $V$ lies in $U$, and $V$ is
open: its complement is the set of $z$ for which some $w$ with $r(w,z)$ lies outside $U$, that is, the
projection to the second factor of the compact set $r\cap\bigl((X\setminus U)\times X\bigr)$, which is
compact, hence closed.
\end{proof}

\formal{\lean{HPC.Jones.isUSCDecomp_of_isClosed}.}

\section{The decomposition space}

The \emph{decomposition space} $X/r$ is the set of elements, with the quotient topology: a set
$\mathcal O$ of elements is open if and only if the union of its members is open in $X$. Let
$\pi\colon X\to X/r$, $\pi(x)=\elt x$, the quotient map; it is continuous, and $X/r=\pi(X)$ is compact.
For $U\subseteq X$ let $U^\sharp$ be the set of elements contained in $U$; if $U$ is open and the
decomposition is u.s.c., then $U^\sharp$ is open in $X/r$, since the union of its members is the open
set $U^*$ of Lemma~\ref{lem:usc-open}.

\begin{theorem}\label{thm:quotient-metrizable}
Let the decomposition be u.s.c. Then:
\begin{enumerate}[label=(\alph*)]
\item the quotient map $\pi$ is closed;
\item $X/r$ is Hausdorff;
\item $X/r$ has a countable basis;
\item $X/r$ is a compact metrizable space.
\end{enumerate}
\end{theorem}

\begin{proof}
(a) Let $F\subseteq X$ be closed. The union of the elements in $\pi(F)$ is the set of points whose element
meets $F$, the complement of $(X\setminus F)^*$, which is closed by Lemma~\ref{lem:usc-open}. So
$\pi(F)$ is closed.

(b) Two distinct elements $\elt x$, $\elt y$ are disjoint closed sets, so they have disjoint open
neighbourhoods $U,V$ in $X$. Then $U^\sharp$ and $V^\sharp$ are disjoint open sets of $X/r$ containing
$\elt x$ and $\elt y$.

(c) Let $\mathcal B$ be a countable basis of $X$, and consider the countably many open sets
$(\bigcup t)^\sharp$, $t$ a finite subfamily of $\mathcal B$. They form a basis: if $\elt x$ lies in an
open set $\mathcal O$ of $X/r$, the open set $\pi^{-1}(\mathcal O)$ contains the compact set $\elt x$, so
finitely many members of $\mathcal B$ contained in $\pi^{-1}(\mathcal O)$ cover $\elt x$; with $t$ the
family of them, $\elt x\in(\bigcup t)^\sharp\subseteq\mathcal O$.

(d) $X/r$ is compact and Hausdorff, hence regular, and has a countable basis; by Urysohn's
metrization theorem it is metrizable.
\end{proof}

\formal{\lean{HPC.Jones.isClosedMap_mk}, \lean{HPC.Jones.t2Space_quotient},
\lean{HPC.Jones.secondCountable_quotient}, \lean{HPC.Jones.metrizableSpace_quotient}. Urysohn's
theorem is Mathlib's \lean{TopologicalSpace.metrizableSpace_of_t3_secondCountable}.}

A decomposition is \emph{monotone} if its elements are connected.

\begin{theorem}\label{thm:monotone-preimage}
Let the decomposition be u.s.c.\ and monotone. If $K\subseteq X/r$ is compact and connected, then
$\pi^{-1}(K)$ is a continuum.
\end{theorem}

\begin{proof}
$K$ is closed ($X/r$ is Hausdorff), so $\pi^{-1}(K)$ is closed, hence compact; it is nonempty. Let
$u,v$ be closed sets with $\pi^{-1}(K)\subseteq u\cup v$ and $\pi^{-1}(K)\cap u\cap v=\emptyset$. Each element
in $K$ is connected and lies in $\pi^{-1}(K)$, so it lies in $u$ or in $v$. The sets
$K_u=\pi(\pi^{-1}(K)\cap u)$ and $K_v=\pi(\pi^{-1}(K)\cap v)$ are closed (Theorem
\ref{thm:quotient-metrizable}(a)) and cover $K$; and $K\cap K_u\cap K_v=\emptyset$, since an element in
both would have a point in $u$ and a point in $v$, while lying entirely in $u$ or entirely in $v$. As
$K$ is connected, $K\subseteq K_u$ or $K\subseteq K_v$; say $K\subseteq K_u$. Then every element in $K$ has a point
in $u$, hence lies in $u$, and $\pi^{-1}(K)\subseteq u$.
\end{proof}

\formal{\lean{HPC.Jones.isContinuum_preimage}.}

\begin{corollary}\label{cor:monotone-decomposable}
Let $X$ be a continuum and the decomposition u.s.c.\ and monotone. If $X/r$ is decomposable, so is
$X$; equivalently, if $X$ is indecomposable, so is $X/r$.
\end{corollary}

\begin{proof}
If $X/r=H\cup K$ with $H,K$ proper subcontinua, then $X=\pi^{-1}(H)\cup\pi^{-1}(K)$, and these are
subcontinua of $X$ by Theorem~\ref{thm:monotone-preimage}, proper because $\pi$ is onto.
\end{proof}

\formal{For monotone maps onto a metric space, \lean{HPC.Hagopian.isDecomposableSpace_of_monotone_image}.}

\begin{proposition}\label{prop:homeo-descends}
Let $T\colon X\to X$ be a homeomorphism permuting the elements: $r(x,y)\iff r(Tx,Ty)$. Then $T$ induces a
homeomorphism $\hat T$ of $X/r$ with $\hat T(\elt x)=\elt{Tx}$.
\end{proposition}

\begin{proof}
$\hat T$ is well defined and bijective, with inverse induced by $T^{-1}$ (which also permutes the
elements). It is continuous because $\hat T\circ\pi=\pi\circ T$ is continuous and $X/r$ carries the
quotient topology; likewise for its inverse.
\end{proof}

\formal{\lean{HPC.Jones.exists_homeomorph_quotient}.}

\section{Continuity at most elements}

A subset of a topological space is \emph{nowhere dense} if its closure has empty interior,
\emph{meager} if it is a countable union of nowhere dense sets, and \emph{residual} if its complement
is meager. In a compact metric space a residual set is dense, by the Baire category theorem
(\vtwo{Lemma~1.8}).

\begin{lemma}\label{lem:frontier-nd}
The frontier $\partial A=A\setminus\interior A$ of a closed set $A$ is closed and nowhere dense.
\end{lemma}

\begin{proof}
It is closed. An open set contained in $A\setminus\interior A$ is contained in $A$, hence in $\interior
A$, hence empty.
\end{proof}

\formal{\lean{HPC.Jones.dense_compl_frontier_of_isClosed}.}

\begin{theorem}[Kuratowski, Fort]\label{thm:residual}
Let the decomposition be u.s.c. Then the set of elements at which it is lower semicontinuous is
residual in $X/r$. In particular (if $X\ne\emptyset$) there is an element at which it is lower
semicontinuous.
\end{theorem}

\begin{proof}
Let $\mathcal B$ be a countable basis of $X$. For $b\in\mathcal B$ the set $A_b=\pi(\cl b)$ is closed
(Theorem~\ref{thm:quotient-metrizable}(a)), so its frontier is closed and nowhere dense, and
$G=\bigcap_{b\in\mathcal B}\bigl(X/r\setminus\partial A_b\bigr)$ is residual. We show that the
decomposition is lower semicontinuous at every element $q=\elt x\in G$.

Let $y\in\elt x$ and $V\ni y$ open, and suppose that every open $W\supseteq\elt x$ contains a point $z$ with
$\elt z\cap V=\emptyset$. Choose $\varepsilon>0$ with $B(y,\varepsilon)\subseteq V$ and $b\in\mathcal B$ with $y\in
b\subseteq B(y,\varepsilon/2)$; then $\cl b\subseteq V$. Now $q\in A_b$ (as $y\in\cl b$) and $q\notin\partial A_b$,
so $q\in\interior A_b$, and $W=\pi^{-1}(\interior A_b)$ is an open set containing $\elt x$. Take $z\in W$
with $\elt z\cap V=\emptyset$. Since $\elt z\in A_b=\pi(\cl b)$, there is $w\in\cl b$ with $\elt w=\elt z$;
then $w\in\elt z\cap V$, a contradiction.

For the last statement: $X/r$ is a nonempty compact metrizable space, in which residual sets are
dense, hence nonempty.
\end{proof}

\formal{\lean{HPC.Jones.residual_lscAt} and \lean{HPC.Jones.exists_lscAt}.}

\begin{proposition}\label{prop:lsc-homeo}
Let $T\colon X\to X$ be a homeomorphism permuting the elements. If the decomposition is lower
semicontinuous at $\elt x$, it is lower semicontinuous at $\elt{Tx}$.
\end{proposition}

\begin{proof}
$T$ maps elements onto elements: $\elt{Ta}=T(\elt a)$. Let $y'\in\elt{Tx}$, so $y'=Ty$ with $y\in\elt x$,
and let $V\ni y'$ be open. Apply lower semicontinuity at $\elt x$ to $y$ and $T^{-1}(V)$: there is an
open $W\supseteq\elt x$ such that every element meeting $W$ meets $T^{-1}(V)$. Then $T(W)$ is an open set
containing $\elt{Tx}$, and an element $\elt{z}$ meeting $T(W)$ is $T(\elt{T^{-1}z})$ with
$\elt{T^{-1}z}$ meeting $W$, hence meeting $T^{-1}(V)$; so $\elt z$ meets $V$.
\end{proof}

\formal{\lean{HPC.Jones.lscAt_homeomorph}.}

\begin{remark}
Theorem~\ref{thm:residual} and Proposition~\ref{prop:lsc-homeo} are used together in Volume~7: if
the homeomorphisms of a homogeneous continuum permute the elements of an u.s.c.\ decomposition,
lower semicontinuity at one element (which exists by Theorem~\ref{thm:residual}) spreads to all
elements, so the decomposition is continuous. This is how F.~B.~Jones's decomposition of a
decomposable homogeneous continuum is shown to be continuous.
\end{remark}

\section*{Exercises}

\begin{exercise}
Show that a decomposition of a compact metric space into closed sets is u.s.c.\ if and only if the
relation $r$ is closed in $X\times X$.
\end{exercise}

\begin{exercise}
Let $f\colon X\to[0,1]$ be continuous. Show that the decomposition of $X$ into the components of
the fibres $f^{-1}(t)$, $t\in[0,1]$, is u.s.c.\ and monotone.
\end{exercise}

\begin{exercise}
Let the decomposition of $[0,1]^2$ consist of the segments $\set t\times[0,1]$ for $t\le\frac12$ and the
points $(t,s)$ for $t>\frac12$. Is it u.s.c.? At which elements is it lower semicontinuous? Check
Theorem~\ref{thm:residual}.
\end{exercise}

%======================================================================
\chapter{The homeomorphism group and the Effros theorem}\label{ch:effros}
%======================================================================

A space $Y$ is \emph{homogeneous} if for any $x,y\in Y$ some homeomorphism $h\colon Y\to Y$ has
$h(x)=y$. The Effros theorem says that for a compact metric $Y$ this can be done
\emph{quantitatively}: points that are close are moved one to the other by homeomorphisms close to
the identity (Theorem~\ref{thm:ungar}). We follow J.~van Mill, \emph{A note on the Effros
theorem} (Amer.\ Math.\ Monthly 111, 2004), as the formal proof does.

This chapter uses some standard notions of descriptive set theory, which we recall without proof
(they are part of Mathlib):
\begin{itemize}
\item A \emph{Polish space} is a separable topological space whose topology comes from a complete
metric. Closed subsets and open subsets of Polish spaces are Polish; a Polish space is
second countable, hence \emph{Lindel\"of}: every open cover has a countable subcover.
\item For a compact metric space $Y$, the space $C(Y,Y)$ of continuous maps with the uniform metric
$d(f,g)=\sup_yd(f(y),g(y))$ is Polish (complete and separable).
\item A set is \emph{analytic} if it is empty or a continuous image of a Polish space.
\item A space is a \emph{Baire space} if every nonempty open subset is non-meager (meager sets were
defined in Chapter~\ref{ch:decomp}). Compact metric spaces are Baire spaces; this is the Baire
category theorem (\vtwo{Lemma~1.8} and its proof).
\end{itemize}

\section{The homeomorphism group}

Let $Y$ be a compact metric space and $\Homeo(Y)$ the set of homeomorphisms $Y\to Y$. It is a group
under composition, $(gh)(y)=g(h(y))$, with identity $1=\mathrm{id}$ and inverse $h^{-1}$, and it acts on
$Y$ by $h\cdot y=h(y)$. We give it the topology induced by the injective map
\[
\iota\colon\Homeo(Y)\to C(Y,Y)\times C(Y,Y),\qquad \iota(h)=(h,h^{-1}),
\]
that is, by the metric $\rho(g,h)=\max\bigl(d(g,h),d(g^{-1},h^{-1})\bigr)$. For $h\in\Homeo(Y)$ let
\[
\mv(h)=\sup_{z\in Y}d(h(z),z)=d(h,\mathrm{id}),
\]
the largest distance a point is moved.

\begin{proposition}\label{prop:homeo-group}
\begin{enumerate}[label=(\alph*)]
\item $\Homeo(Y)$ is a topological group, and the action $\Homeo(Y)\times Y\to Y$ is continuous.
\item $\Homeo(Y)$ is a Polish space.
\item $\mv(h^{-1})=\mv(h)$ and $\mv(gh)\le\mv(g)+\mv(h)$; and the sets $\set{h:\mv(h)<r}$, $r>0$,
form a basis of neighbourhoods of $1$.
\end{enumerate}
\end{proposition}

\begin{proof}
(a) Composition $C(Y,Y)\times C(Y,Y)\to C(Y,Y)$ is continuous: $d(f\circ g,f'\circ g')\le d(f\circ g,f\circ
g')+d(f,f')$, and the first term is small when $g'$ is uniformly close to $g$, by uniform continuity of
$f$. Since $\iota(gh)=(g\circ h,\ h^{-1}\circ g^{-1})$, multiplication is continuous; inversion is
continuous since $\iota(h^{-1})$ is $\iota(h)$ with its coordinates swapped. Evaluation is continuous:
$d(h(y),h'(y'))\le d(h(y),h(y'))+d(h,h')$.

(b) The image of $\iota$ is $\set{(f,g):f\circ g=\mathrm{id},\ g\circ f=\mathrm{id}}$ (a pair of mutually
inverse continuous maps is a homeomorphism and its inverse), which is closed in the Polish space
$C(Y,Y)^2$ because composition is continuous. So $\iota$ is a homeomorphism onto a closed subset of a
Polish space, and $\Homeo(Y)$ is Polish.

(c) $\mv(h^{-1})=\sup_zd(h^{-1}(z),z)=\sup_wd(w,h(w))=\mv(h)$ (substitute $z=h(w)$), and
$d(g(h(z)),z)\le d(g(h(z)),h(z))+d(h(z),z)$. The $\rho$-ball of radius $r$ about $1$ is
$\set{h:\mv(h)<r,\ \mv(h^{-1})<r}=\set{h:\mv(h)<r}$.
\end{proof}

\formal{\lean{HPC.Effros.homeoGroup}, \lean{HPC.Effros.homeoTopology},
\lean{HPC.Effros.homeo_isTopologicalGroup}, \lean{HPC.Effros.homeo_continuousSMul},
\lean{HPC.Effros.homeo_polishSpace}, \lean{HPC.Effros.supMove} (our $\mv$),
\lean{HPC.Effros.supMove_inv}, \lean{HPC.Effros.supMove_mul_le}, \lean{HPC.Effros.hasBasis_supMove}
(in \lean{HPC/Effros/HomeoGroup.lean}). Mathlib has no topology on the homeomorphism group at the
pinned version; the formal proof uses the compact-open topology on $C(Y,Y)$, which for compact
metric $Y$ is the topology of the uniform metric.}

\section{Meager, nowhere meager and analytic sets}

In this section $Y$ is a topological space.

\begin{definition}\label{def:fat}
A set $A\subseteq Y$ is \emph{nowhere meager} if for every open $V$ with $V\cap A\neq\emptyset$ the set
$V\cap A$ is not meager in $Y$. It is \emph{fat} if it is dense and nowhere meager.
\end{definition}

\formal{\lean{HPC.Effros.NowhereMeagre}, \lean{HPC.Effros.IsFat}.}

\begin{lemma}\label{lem:fat-basic}
\begin{enumerate}[label=(\alph*)]
\item If $A$ is nowhere meager and $M$ is meager, then $A\setminus M$ is nowhere meager, and
$\cl A\subseteq\cl{A\setminus M}$.
\item A set that is not meager has a closure with nonempty interior.
\end{enumerate}
\end{lemma}

\begin{proof}
(a) If $V$ is open and $V\cap A\ne\emptyset$, then $V\cap A\subseteq(V\cap(A\setminus M))\cup M$; the left side
is not meager, so $V\cap(A\setminus M)$ is not meager, in particular nonempty. Taking for $V$ the
neighbourhoods of a point of $\cl A$ gives the second statement.
(b) Otherwise the set is nowhere dense, hence meager.
\end{proof}

\formal{\lean{HPC.Effros.NowhereMeagre.diff_meagre} (van Mill's Lemma~2.1),
\lean{HPC.Effros.NowhereMeagre.closure_subset_closure_diff},
\lean{HPC.Effros.nonempty_interior_closure_of_not_meagre}.}

The key fact about analytic sets is the following intersection property.

\begin{lemma}[The nested construction]\label{lem:core-construction}
Let $P,Q$ be complete metric spaces with $P\ne\emptyset$, $Y$ a metric space, and $\alpha\colon P\to Y$,
$\beta\colon Q\to Y$ continuous maps such that
\begin{enumerate}[label=(\roman*)]
\item for every nonempty open $U\subseteq P$, $\interior\cl{\alpha(U)}\neq\emptyset$, and for every nonempty
open $V\subseteq Q$, $\interior\cl{\beta(V)}\ne\emptyset$;
\item $\alpha(P)\subseteq\cl{\beta(Q)}$.
\end{enumerate}
Then $\alpha(p)=\beta(q)$ for some $p\in P$, $q\in Q$.
\end{lemma}

\begin{proof}
Call a triple $(U,V,W)$ of open sets of $P$, $Q$, $Y$ \emph{good} if $U\ne\emptyset$, $\beta(V)\subseteq W$
and $\alpha(U)\subseteq\cl{\beta(V)}$. The triple $(P,Q,Y)$ is good by (ii).

\emph{Step.} Let $(U,V,W)$ be good and $\varepsilon>0$. We find a good triple $(E,F,W')$ with $E\subseteq U$,
$F\subseteq V$, $W'\subseteq W$, each contained in a ball of radius $\varepsilon$. Let
$O=\interior\cl{\alpha(U)}$, nonempty by (i). Since $\alpha(U)\subseteq\cl{\beta(V)}\subseteq\cl W$, the
nonempty open set $O$ lies in $\cl W$, so it meets $W$: take $y\in O\cap W$ and put $W'=O\cap W\cap
B(y,\varepsilon)$. Then $y\in W'\subseteq\cl{\alpha(U)}\subseteq\cl{\beta(V)}$, so the open set $W'$ meets
$\beta(V)$: take $q\in V$ with $\beta(q)\in W'$, and put $F=V\cap\beta^{-1}(W')\cap B(q,\varepsilon)$, a
nonempty open set with $\beta(F)\subseteq W'$. Let $O'=\interior\cl{\beta(F)}$, nonempty by (i); it lies in
$\cl{W'}\subseteq\cl{\alpha(U)}$, so it meets $\alpha(U)$: take $p\in U$ with $\alpha(p)\in O'$, and put
$E=U\cap\alpha^{-1}(O')\cap B(p,\varepsilon)$, nonempty and open, with $\alpha(E)\subseteq
O'\subseteq\cl{\beta(F)}$. The triple $(E,F,W')$ is good.

\emph{Limit.} Starting from $(P,Q,Y)$ and applying the step with $\varepsilon=2^{-n}$ at stage $n$, we
get good triples $(U_n,V_n,W_n)$, decreasing in each coordinate, with $U_{n+1},V_{n+1},W_{n+1}$ in balls
of radius $2^{-n}$. Each $V_n$ is nonempty, since $\emptyset\ne\alpha(U_n)\subseteq\cl{\beta(V_n)}$. Choosing a
point in each $U_n$ gives a Cauchy sequence; by completeness it converges to some $p$, which lies in
every $\cl{U_n}$. Likewise there is $q$ in every $\cl{V_n}$. Then $\alpha(p)\in\cl{\alpha(U_n)}\subseteq\cl{W_n}$
and $\beta(q)\in\cl{\beta(V_n)}\subseteq\cl{W_n}$ for every $n$, and $\diam\cl{W_{n+1}}\le2\cdot2^{-n}\to0$.
So $\alpha(p)=\beta(q)$.
\end{proof}

\formal{\lean{HPC.Effros.CoreGood}, \lean{HPC.Effros.coreGood_step},
\lean{HPC.Effros.exists_eq_of_core}.}

\begin{lemma}[Reducing the parameter space]\label{lem:reduced}
Let $P$ be a Polish space, $Y$ a topological space and $f\colon P\to Y$ any map. There is a closed set
$C\subseteq P$ such that $f(P)\setminus f(C)$ is meager, and $f(U)$ is not meager for every nonempty
relatively open subset $U$ of $C$.
\end{lemma}

\begin{proof}
Let $O$ be the union of all open sets $U'\subseteq P$ with $f(U')$ meager. Each point of $O$ lies in a
basic open set (of a fixed countable basis) contained in such a $U'$, whose image is then meager;
so $O$ is a countable union of open sets with meager images, and $f(O)$ is meager. Let $C=P\setminus
O$. Then $f(P)\setminus f(C)\subseteq f(O)$ is meager. If $U=C\cap U'$, with $U'$ open, is nonempty and
$f(U)$ is meager, then $f(U')\subseteq f(U)\cup f(O)$ is meager, so $U'\subseteq O$ and $U=\emptyset$, a
contradiction.
\end{proof}

\formal{\lean{HPC.Effros.exists_reduced_param} (stated for $P=\N^\N$, which suffices since Mathlib
defines analytic sets as continuous images of $\N^\N$).}

\begin{theorem}[van Mill, Proposition~2.2, relative form]\label{thm:prop22}
Let $Y$ be a metric space and $S,T\subseteq Y$ analytic and nowhere meager, with $S\neq\emptyset$ and
$S\subseteq\cl T$. Then $S\cap T\neq\emptyset$.
\end{theorem}

\begin{proof}
$T$ is nonempty since $S\subseteq\cl T$. Write $S=f(P)$ and $T=g(Q)$ with $P,Q$ Polish and $f,g$ continuous,
and let $C\subseteq P$ and $D\subseteq Q$ be given by Lemma~\ref{lem:reduced}. They are closed, hence Polish;
fix complete metrics on them. $S$ is not meager (apply the definition of nowhere meager with
$V=Y$), and $S\subseteq f(C)\cup(f(P)\setminus f(C))$, so $f(C)\ne\emptyset$ and $C\ne\emptyset$.

Apply Lemma~\ref{lem:core-construction} to $\alpha=f|_C$ and $\beta=g|_D$. Condition (i) holds by
Lemma~\ref{lem:reduced} and Lemma~\ref{lem:fat-basic}(b). For (ii), let $M=g(Q)\setminus g(D)$, a meager
set; then $T\setminus M\subseteq g(D)$, and by Lemma~\ref{lem:fat-basic}(a)
\[
f(C)\subseteq S\subseteq\cl T\subseteq\cl{T\setminus M}\subseteq\cl{g(D)} .
\]
So $f(p)=g(q)$ for some $p\in C$, $q\in D$, a point of $S\cap T$.
\end{proof}

\formal{\lean{HPC.Effros.inter_nonempty_of_nowhereMeagre}.}

\begin{corollary}[van Mill, Proposition~2.2]\label{cor:fat-meet}
In a nonempty metric space, two analytic fat sets meet.
\end{corollary}

\formal{\lean{HPC.Effros.inter_nonempty_of_isFat}. The hypothesis that the space be nonempty is
needed: in the empty space the empty set is analytic and fat
(\lean{HPC.Effros.isFat_empty_of_isEmpty}).}

\section{The open mapping principle}

In this section $G$ is a Polish group (a topological group whose topology is Polish), acting
continuously and transitively on a nonempty metric space $X$ which is a Baire space. For $U\subseteq G$
and $x\in X$ let $U\cdot x=\set{g\cdot x:g\in U}$. For each $g\in G$ the map $x\mapsto g\cdot x$ is a
homeomorphism of $X$ (with inverse $x\mapsto g^{-1}\cdot x$); it maps meager sets to meager sets.

\begin{lemma}\label{lem:vm-basis}
There are open sets $G=U_0\supseteq U_1\supseteq\cdots$ containing $1$, with $U_n^{-1}=U_n$ and
$U_{n+1}U_{n+1}\subseteq U_n$, forming a basis of neighbourhoods of $1$.
\end{lemma}

\begin{proof}
$G$ is metrizable, so $1$ has a decreasing basis of neighbourhoods $B_0\supseteq B_1\supseteq\cdots$. Given an
open symmetric $U_n\ni1$, continuity of multiplication at $(1,1)$ gives an open $S\ni1$ with $SS\subseteq
U_n\cap B_n$; put $U_{n+1}=S\cap S^{-1}$.
\end{proof}

\formal{\lean{HPC.Effros.exists_vanMill_basis}.}

Fix such sets $U_n$.

\begin{lemma}[van Mill, Lemma~3.1]\label{lem:vm31}
If $V\subseteq X$ is open and $z\in V\cap U_n\cdot x$, then $U_m\cdot z\subseteq V\cap U_n\cdot x$ for some $m$.
\end{lemma}

\begin{proof}
Write $z=h\cdot x$ with $h\in U_n$. The set $\set{g:(gh)\cdot x\in V,\ gh\in U_n}$ is open (by continuity)
and contains $1$, so it contains some $U_m$. For $g\in U_m$, $g\cdot z=(gh)\cdot x\in V$ with $gh\in U_n$.
\end{proof}

\formal{\lean{HPC.Effros.vanMill_3_1}.}

\begin{lemma}[van Mill, Lemma~3.2]\label{lem:vm32}
For every neighbourhood $U$ of $1$ and every $x$, the set $U\cdot x$ is not meager.
\end{lemma}

\begin{proof}
Let $U'=\interior U\ni1$. The open sets $\varphi U'$, $\varphi\in G$, cover $G$; as $G$ is Lindel\"of,
countably many of them, $\varphi_kU'$, cover $G$. Every $y\in X$ is $g\cdot x$ for some $g$ (transitivity),
and $g=\varphi_ku$ with $u\in U'$, so $y=\varphi_k\cdot(u\cdot x)$. Hence $X=\bigcup_k\varphi_k\cdot(U\cdot x)$.
If $U\cdot x$ were meager, so would be each $\varphi_k\cdot(U\cdot x)$, and $X$ would be meager, which is
impossible in a nonempty Baire space.
\end{proof}

\formal{\lean{HPC.Effros.vanMill_3_2}.}

\begin{lemma}[van Mill, Corollary~3.3]\label{lem:vm33}
Each $U_n\cdot x$ is nowhere meager.
\end{lemma}

\begin{proof}
If $V$ is open and $z\in V\cap U_n\cdot x$, Lemma~\ref{lem:vm31} gives $U_m\cdot z\subseteq V\cap U_n\cdot x$,
and $U_m\cdot z$ is not meager by Lemma~\ref{lem:vm32}.
\end{proof}

\formal{\lean{HPC.Effros.vanMill_3_3}.}

\begin{lemma}[van Mill, Lemma~3.4]\label{lem:vm34}
$x\in\interior\cl{U_n\cdot x}$ for every $n$ and $x$.
\end{lemma}

\begin{proof}
$V=\interior\cl{U_{n+1}\cdot x}$ is nonempty by Lemmas~\ref{lem:vm32} and~\ref{lem:fat-basic}(b), and,
being a nonempty open subset of $\cl{U_{n+1}\cdot x}$, it meets $U_{n+1}\cdot x$: say $h\cdot x\in V$ with
$h\in U_{n+1}$. The set $h^{-1}\cdot V$ is open, contains $x$, and
\[
h^{-1}\cdot V\subseteq h^{-1}\cdot\cl{U_{n+1}\cdot x}=\cl{(h^{-1}U_{n+1})\cdot x}\subseteq\cl{U_n\cdot x},
\]
because $h^{-1}\in U_{n+1}$ and $U_{n+1}U_{n+1}\subseteq U_n$.
\end{proof}

\formal{\lean{HPC.Effros.vanMill_3_4} (as printed by van Mill, with the closure).}

\begin{proposition}[van Mill, Proposition~3.5]\label{prop:vm35}
$\interior\cl{U_{n+1}\cdot x}\subseteq U_n\cdot x$ for every $n$ and $x$.
\end{proposition}

\begin{proof}
Let $z\in\interior\cl{U_{n+1}\cdot x}$, and let $E=\interior\cl{U_{n+1}\cdot x}\cap\interior\cl{U_{n+1}\cdot
z}$, an open set containing $z$ (Lemma~\ref{lem:vm34}). Consider
\[
S'=U_{n+1}\cdot x\cap E,\qquad T'=U_{n+1}\cdot z\cap E .
\]
\begin{itemize}
\item They are analytic: $T'$ is the image of the open, hence Polish, set
$\set{g\in U_{n+1}:g\cdot z\in E}$ under the continuous map $g\mapsto g\cdot z$; similarly for $S'$.
\item They are nowhere meager: for $V$ open, $V\cap T'=(V\cap E)\cap U_{n+1}\cdot z$, which is not meager
if nonempty, by Lemma~\ref{lem:vm33}; similarly for $S'$.
\item $E\subseteq\cl{T'}$: if $y\in E$ and $N$ is a neighbourhood of $y$, then $N\cap E$ is a neighbourhood of
$y\in\cl{U_{n+1}\cdot z}$, so it meets $U_{n+1}\cdot z$, that is, $N$ meets $T'$. Likewise $E\subseteq\cl{S'}$;
in particular $S'\ne\emptyset$ (as $z\in E$), and $S'\subseteq E\subseteq\cl{T'}$.
\end{itemize}
By Theorem~\ref{thm:prop22}, $S'\cap T'\ne\emptyset$: $g\cdot x=h\cdot z$ with $g,h\in U_{n+1}$. Then
$z=(h^{-1}g)\cdot x$ and $h^{-1}g\in U_{n+1}U_{n+1}\subseteq U_n$.
\end{proof}

\formal{\lean{HPC.Effros.vanMill_3_5} (as printed, with the closure).}

\begin{theorem}[The open mapping principle; Effros]\label{thm:effros-open}
Let a Polish group $G$ act continuously and transitively on a nonempty metric Baire space $X$. Then
for every $x\in X$ and every neighbourhood $U$ of $1$ in $G$, the set $U\cdot x$ is a neighbourhood of
$x$. Consequently the orbit maps $g\mapsto g\cdot x$ are open.
\end{theorem}

\begin{proof}
Choose $n$ with $U_n\subseteq U$. By Lemma~\ref{lem:vm34} and Proposition~\ref{prop:vm35},
$x\in\interior\cl{U_{n+1}\cdot x}\subseteq U_n\cdot x\subseteq U\cdot x$. For the second statement, let $O\subseteq G$
be open and $g\in O$; then $g^{-1}O$ is a neighbourhood of $1$, so $(g^{-1}O)\cdot x$ is a neighbourhood
of $x$, and its image $O\cdot x$ under the homeomorphism $y\mapsto g\cdot y$ is a neighbourhood of $g\cdot
x$.
\end{proof}

\formal{\lean{HPC.Effros.effros_open_mapping} and \lean{HPC.Effros.isOpenMap_smul_left}.}

\section{Ungar's corollary}

\begin{theorem}[Effros, Ungar]\label{thm:ungar}
Let $Y$ be a homogeneous compact metric space and $\varepsilon>0$. There is $\delta>0$ such that for any
$x,y\in Y$ with $d(x,y)<\delta$ there is a homeomorphism $h\colon Y\to Y$ with $h(x)=y$ and
$d(h(z),z)<\varepsilon$ for every $z\in Y$.
\end{theorem}

\begin{proof}
If $Y=\emptyset$ there is nothing to prove. By Proposition~\ref{prop:homeo-group}, $G=\Homeo(Y)$ is a
Polish group acting continuously on $Y$; the action is transitive because $Y$ is homogeneous, and $Y$
is a nonempty compact metric space, hence a Baire space. Let $N=\set{h:\mv(h)<\varepsilon/2}$, a
neighbourhood of $1$. By Theorem~\ref{thm:effros-open} each $N\cdot z$ is a neighbourhood of $z$, so the
interiors of the sets $N\cdot z$, $z\in Y$, form an open cover of the compact space $Y$. Let $\delta$ be
a Lebesgue number: every ball $B(x,\delta)$ lies in some $\interior(N\cdot z)$.

Let $d(x,y)<\delta$, and take $z$ with $x,y\in N\cdot z$: $x=g(z)$, $y=h(z)$ with $g,h\in N$. Then
$hg^{-1}$ maps $x$ to $y$, and $\mv(hg^{-1})\le\mv(h)+\mv(g^{-1})=\mv(h)+\mv(g)<\varepsilon$.
\end{proof}

\formal{\lean{HPC.Effros.effros_homogeneous}; for a compact subset $M$ of a metric space which is
homogeneous in its own topology (the form used for plane continua),
\lean{HPC.Effros.effros_homogeneous_set}.}

\begin{remark}
In the later volumes Theorem~\ref{thm:ungar} is used to move small subcontinua of a homogeneous
continuum around by homeomorphisms that hardly move anything, which is the basic tool for turning
homogeneity into local structure; the formal proof uses it in its treatment of F.~B.~Jones's
papers of 1951, 1955 and 1975, of the span argument of Oversteegen and Tymchatyn, and of Hagopian's
1984 paper.
\end{remark}

\section*{Exercises}

\begin{exercise}
Show that the conclusion of Theorem~\ref{thm:ungar} fails for the non-homogeneous space $[0,1]$: there
is no homeomorphism moving $0$ to a point $y>0$.
\end{exercise}

\begin{exercise}
Prove Theorem~\ref{thm:ungar} directly for the circle, with an explicit $\delta$.
\end{exercise}

\begin{exercise}
Let $Y$ be a homogeneous compact metric space. Show that for every $x\in Y$ and $\varepsilon>0$ the set
of points $y$ for which some homeomorphism $h$ with $\mv(h)<\varepsilon$ maps $x$ to $y$ is open.
\end{exercise}


\appendix
%======================================================================
\chapter{Correspondence with the formal proof}\label{app:formal}
%======================================================================

The table lists the definitions and results of this volume that carry a ``Formal counterpart''
line, together with the corresponding declarations of the Lean development. All declaration
names start with \lean{HPC.}, except the few Mathlib names marked as such; file names are relative
to the root of the repository. The formal statements are sometimes slightly more general, or are
split differently into lemmas; the table points to the place where the content of each result is
proved. Rows are in the order of the text.

The main theorems of the volume are, in the Lean development (prefix \lean{HPC.} omitted):
\begin{center}
\begin{tabular}{@{}p{0.42\textwidth}p{0.54\textwidth}@{}}
\toprule
Boundary bumping (Theorem~\ref{thm:bump-closed}) & \lean{Jones.exists_continuum_bump_compact},
\lean{Plane.exists_continuum_bump_closed}\\
Moore's non-cut point theorem (Theorems~\ref{thm:noncut-side}, \ref{thm:two-noncut}) &
\lean{Plane.exists_nonCutPoint_of_cut}, \lean{Plane.exists_two_nonCutPoints}\\
Sorgenfrey's theorem (Theorem~\ref{thm:sorgenfrey}) & \lean{Plane.sorgenfrey}\\
Snake-like continua are hereditarily unicoherent and atriodic
(Theorems~\ref{thm:snake-unicoherent}, \ref{thm:snake-atriodic}) &
\lean{Plane.isHereditarilyUnicoherent_of_isSnakeLike}, \lean{Plane.isAtriodic_of_isSnakeLike}\\
Snake-likeness is topological (Theorem~\ref{thm:snake-homeo}) & \lean{BJ59.isSnakeLike_of_homeomorph}\\
Decomposition spaces are metrizable (Theorem~\ref{thm:quotient-metrizable}) &
\lean{Jones.metrizableSpace_quotient}\\
Continuity at a residual set of elements (Theorem~\ref{thm:residual}) & \lean{Jones.residual_lscAt}\\
The open mapping principle (Theorem~\ref{thm:effros-open}) & \lean{Effros.effros_open_mapping}\\
Ungar's corollary (Theorem~\ref{thm:ungar}) & \lean{Effros.effros_homogeneous}\\
\bottomrule
\end{tabular}
\end{center}
All of them are proved in the Lean development without \lean{sorry} and without axioms beyond the
standard ones of Lean and Mathlib.

\begin{longtable}{@{}p{0.40\textwidth}p{0.56\textwidth}@{}}
\toprule
In these notes & In the Lean development\\
\midrule
\endhead
\bottomrule
\endfoot
\multicolumn{2}{@{}l}{\emph{Chapter~\ref{ch:bumping}}}\\
Lemma~\ref{lem:sep-of-not-conn} & \lean{HPC.Plane.exists_separated_of_not_isPreconnected}.\\
Lemma~\ref{lem:side} & \lean{HPC.Plane.subset_side_of_isPreconnected}.\\
Lemma~\ref{lem:union-continua} & \lean{HPC.Plane.isContinuum_union}.\\
A continuum plus one piece (Lemma~\ref{lem:piece}) & \lean{HPC.Plane.isContinuum_union_piece} (Nadler, \emph{Continuum Theory}, 6.3). A version for an arbitrary connected topological space is Theorem~\ref{thm:nadler63}.\\
Lemma~\ref{lem:nested} & \lean{HPC.T14.isPreconnected_sInter_chain}.\\
Lemma~\ref{lem:quasi} & Mathlib's \lean{connectedComponent_eq_iInter_isClopen}, used inside the proofs below.\\
Lemma~\ref{lem:clopen-split} & \lean{HPC.Plane.exists_clopen_split}.\\
Separating two closed sets in a compact set (Theorem~\ref{thm:closed-split}) & \lean{HPC.Plane.exists_closed_split_of_no_connected} (Moore's Theorem~44 in \emph{Foundations of Point Set Theory}). A variant with a continuum in place of the preconnected set, in a compact space, is \lean{HPC.Plane.closed_dichotomy}.\\
Boundary bumping (Theorem~\ref{thm:bump-closed}) & \lean{HPC.Jones.exists_continuum_bump_compact}; for $M$ a continuum in a compact space and $F$ closed, \lean{HPC.Plane.exists_continuum_bump_closed} (Nadler 5.4).\\
Boundary bumping, open form (Corollary~\ref{cor:bump-open}) & \lean{HPC.Plane.exists_continuum_bump_open} (Nadler 5.7).\\
Boundary bumping, metric form (Corollary~\ref{cor:bump-ball}) & \lean{HPC.Plane.exists_subcontinuum_meeting_closedBall}.\\[4pt]
\multicolumn{2}{@{}l}{\emph{Chapter~\ref{ch:irreducible}}}\\
Definition~\ref{def:irreducible} & \lean{HPC.IsIrreducibleBetween} (in \lean{HPC/PseudoArc/Theorem14.lean}).\\
Definition~\ref{def:irr-sets} & \lean{HPC.Plane.IrrFromTo} (Nadler 11.29).\\
Existence of irreducible subcontinua (Theorem~\ref{thm:irr-exists}) & \lean{HPC.Plane.exists_irrFromTo} (Nadler 11.30) and \lean{HPC.exists_irreducible}.\\
Theorem~\ref{thm:indec-irr} & \lean{HPC.Plane.exists_irreducible_of_indecomposable} (the case of Nadler 11.15 used later).\\
Proposition~\ref{prop:irr-homeo} & \lean{HPC.IsIrreducibleBetween.homImage}.\\[4pt]
\multicolumn{2}{@{}l}{\emph{Chapter~\ref{ch:noncut}}}\\
Definition~\ref{def:separation} & \lean{HPC.Plane.IsSeparation}, \lean{HPC.Plane.IsCutPoint} (in \lean{HPC/Plane/NonCut51Gen.lean}).\\
Lemma~\ref{lem:sep-basic} & \lean{HPC.Plane.IsSeparation.disjoint_closure_left}, \lean{HPC.Plane.IsSeparation.isOpen_right}, \lean{HPC.Plane.IsSeparation.subset_or_subset}, \lean{HPC.Plane.exists_separation_of_not_isConnected}.\\
Theorem~\ref{thm:nadler63} & \lean{HPC.Plane.isConnected_union_of_separation} (Nadler 6.3).\\
Lemma~\ref{lem:nadler64} & \lean{HPC.Plane.nadler_6_4}.\\
Moore (Theorem~\ref{thm:noncut-side}) & \lean{HPC.Plane.exists_nonCutPoint_left} and \lean{HPC.Plane.exists_nonCutPoint_of_cut} (Nadler 6.6). Nadler assumes $X$ nondegenerate; this follows from the separation and is not needed.\\
Theorem~\ref{thm:two-noncut} & \lean{HPC.Plane.exists_two_nonCutPoints}.\\
Theorem~\ref{thm:noncut-not-in} & \lean{HPC.Plane.not_subset_of_nonCutPoints} (Nadler 6.7).\\[4pt]
\multicolumn{2}{@{}l}{\emph{Chapter~\ref{ch:triods}}}\\
Definition~\ref{def:triod} & \lean{HPC.Plane.IsTriod}, \lean{HPC.Plane.IsAtriodic} (in \lean{HPC/Plane/Atriodic51Defs.lean}); for a triod in a non-compact space, \lean{HPC.Hagopian.IsTriodSet}.\\
Definition~\ref{def:unicoherent} & \lean{HPC.Plane.IsUnicoherent}, \lean{HPC.Plane.IsHereditarilyUnicoherent}.\\
Definition~\ref{def:esum} & \leansub{HPC.Plane.IsEssentialSum}{2}, \leansub{HPC.Plane.IsEssentialSum}{3}, \lean{HPC.Plane.IsWeakTriod} (in \lean{HPC/Plane/SorgenfreyDefs.lean}); the constructor used in all proofs below is \lean{HPC.Plane.isWeakTriod_of}.\\
Lemma~\ref{lem:esum-basic} & The lemmas in the namespace \leansub{HPC.Plane.IsEssentialSum}{2} (\lean{diff_nonempty}, \lean{closure_diff_subset_left}, \lean{eq_closure_union}, \lean{closure_diff_ne}, \lean{closure_diff_disjoint}, \lean{closure_inter_nonempty}).\\
A triod with a given core (Lemma~\ref{lem:core}) & \lean{HPC.Plane.isTriod_of_core}.\\
Proposition~\ref{prop:triod-weak} & \lean{HPC.Plane.isWeakTriod_of_isTriod} (Nadler 11.25).\\
Corollary~\ref{cor:indec-not-triod} & \lean{HPC.Hagopian.not_isTriod_univ_of_indecomposable} (for the whole space, which is all that is used).\\
Lemma~\ref{lem:weak-not-irr} & \lean{HPC.Plane.not_irreducible_of_isWeakTriod}.\\
Proposition~\ref{prop:uni-triod} & \lean{HPC.Plane.isTriod_iff_isWeakTriod} (Nadler 11.26).\\
Theorem~\ref{thm:diff-connected} & \lean{HPC.Plane.isConnected_diff_of_not_weakTriod} (Nadler 11.27), with the two halves in \lean{HPC.Plane.isConnected_diff_left_of_not_weakTriod}; the continuum $\cl U$ is \lean{HPC.Plane.isContinuum_closure_diff}.\\
Lemma~\ref{lem:1128} & \lean{HPC.Plane.isConnected_diff_of_superset'} (Nadler 11.28, corrected).\\
Proposition~\ref{prop:1131} & \lean{HPC.Plane.nadler_11_31} (Nadler 11.31, without the unused hypothesis $B\subseteq\cl U$).\\
Proposition~\ref{prop:1132} & \lean{HPC.Plane.nadler_11_32} (Nadler 11.32, without the hypothesis $A\ne\cl U$, which follows from the others).\\
Lemma~\ref{lem:1133a} & \lean{HPC.Plane.diff_subset_of_irr}.\\
Proposition~\ref{prop:1133} & \lean{HPC.Plane.nadler_11_33} (Nadler 11.33).\\
Lemma~\ref{lem:sorg-proper} & \lean{HPC.Plane.exists_proper_of_not_irr}.\\
Lemma~\ref{lem:sorg-step} & \lean{HPC.Plane.SorgGood}, \lean{HPC.Plane.sorgGood_base}, \lean{HPC.Plane.sorgGood_step}.\\
Lemma~\ref{lem:sorg-star} & \lean{HPC.Plane.exists_irrFromTo_point} (Nadler 11.34(a)); the body of the proof by contradiction is \lean{HPC.Plane.sorgenfrey_contra}.\\
Sorgenfrey (Theorem~\ref{thm:sorgenfrey}) & \lean{HPC.Plane.sorgenfrey} (Nadler 11.34).\\
Corollary~\ref{cor:irr-iff} & \lean{HPC.Plane.irreducible_iff_not_isTriod} (Nadler 11.35).\\
Corollary~\ref{cor:atriodic-irr} & \lean{HPC.Plane.exists_irreducibleBetween_of_atriodic}.\\[4pt]
\multicolumn{2}{@{}l}{\emph{Chapter~\ref{ch:snake}}}\\
Theorem~\ref{thm:snake-sub} & \lean{HPC.Plane.isSnakeLike_of_subcontinuum}; the interval of links meeting $K$ is \lean{HPC.Plane.meets_between}.\\
Intermediate values along a chain (Lemma~\ref{lem:chain-ivt}) & \lean{HPC.Plane.chain_ivt}.\\
Corollary~\ref{cor:two-sets} & \lean{HPC.Plane.chain_two_sets}.\\
Theorem~\ref{thm:snake-unicoherent} & \lean{HPC.Plane.isUnicoherent_of_isSnakeLike} and \lean{HPC.Plane.isHereditarilyUnicoherent_of_isSnakeLike} (Nadler 12.2).\\
Theorem~\ref{thm:snake-atriodic} & \lean{HPC.Plane.isAtriodic_of_isSnakeLike}, through \lean{HPC.Plane.not_isTriod_of_isSnakeLike} (Nadler 12.4).\\
Proposition~\ref{prop:chain-epsmap} & \lean{HPC.BJ59.exists_epsMap_of_chain} and \lean{HPC.BJ59.exists_epsMap_of_isSnakeLike}.\\
Lemma~\ref{lem:uniform-fibres} & \lean{HPC.Plane.exists_delta_of_fibres_on}.\\
The $\varepsilon$-map criterion (Theorem~\ref{thm:epsmap-criterion}) & \lean{HPC.Plane.isSnakeLike_of_epsMap}; for the whole space, \lean{HPC.Plane.isSnakeLike_univ_of_epsMap}.\\
Snake-likeness is topological (Theorem~\ref{thm:snake-homeo}) & \lean{HPC.BJ59.isSnakeLike_of_homeomorph}.\\
Definition~\ref{def:circchain} & \lean{HPC.Plane.CircChain}, \lean{HPC.Plane.IsCircEpsChain}, \lean{HPC.Plane.IsCircleLike} (in \lean{HPC/Plane/Ingram68.lean}).\\
Proposition~\ref{prop:circle-circlelike} & \lean{HPC.BJ59.isCircleLike_circle}, with the arcs \lean{HPC.BJ59.circArc}.\\
Proposition~\ref{prop:circle-homeo} & \lean{HPC.BJ59.isCircleLike_univ_of_homeomorph}.\\
Theorem~\ref{thm:circle-sub} & \lean{HPC.Plane.isSnakeLike_of_subset_circleLike} (Nadler 12.51).\\[4pt]
\multicolumn{2}{@{}l}{\emph{Chapter~\ref{ch:decomp}}}\\
Definition~\ref{def:usc} & \lean{HPC.Jones.IsUSCDecomp}, \lean{HPC.Jones.LSCAt}, \lean{HPC.Jones.IsContinuousDecomp}, with the element $\elt x$ written \lean{HPC.Jones.cls} (in \lean{HPC/Jones/Decomp.lean}).\\
Lemma~\ref{lem:usc-open} & \lean{HPC.Jones.isOpen_setOf_cls_subset}.\\
Proposition~\ref{prop:closed-usc} & \lean{HPC.Jones.isUSCDecomp_of_isClosed}.\\
Theorem~\ref{thm:quotient-metrizable} & \lean{HPC.Jones.isClosedMap_mk}, \lean{HPC.Jones.t2Space_quotient}, \lean{HPC.Jones.secondCountable_quotient}, \lean{HPC.Jones.metrizableSpace_quotient}. Urysohn's theorem is Mathlib's \lean{TopologicalSpace.metrizableSpace_of_t3_secondCountable}.\\
Theorem~\ref{thm:monotone-preimage} & \lean{HPC.Jones.isContinuum_preimage}.\\
Corollary~\ref{cor:monotone-decomposable} & For monotone maps onto a metric space, \lean{HPC.Hagopian.isDecomposableSpace_of_monotone_image}.\\
Proposition~\ref{prop:homeo-descends} & \lean{HPC.Jones.exists_homeomorph_quotient}.\\
Lemma~\ref{lem:frontier-nd} & \lean{HPC.Jones.dense_compl_frontier_of_isClosed}.\\
Kuratowski, Fort (Theorem~\ref{thm:residual}) & \lean{HPC.Jones.residual_lscAt} and \lean{HPC.Jones.exists_lscAt}.\\
Proposition~\ref{prop:lsc-homeo} & \lean{HPC.Jones.lscAt_homeomorph}.\\[4pt]
\multicolumn{2}{@{}l}{\emph{Chapter~\ref{ch:effros}}}\\
Proposition~\ref{prop:homeo-group} & \lean{HPC.Effros.homeoGroup}, \lean{HPC.Effros.homeoTopology}, \lean{HPC.Effros.homeo_isTopologicalGroup}, \lean{HPC.Effros.homeo_continuousSMul}, \lean{HPC.Effros.homeo_polishSpace}, \lean{HPC.Effros.supMove} (our $\mv$), \lean{HPC.Effros.supMove_inv}, \lean{HPC.Effros.supMove_mul_le}, \lean{HPC.Effros.hasBasis_supMove} (in \lean{HPC/Effros/HomeoGroup.lean}). Mathlib has no topology on the homeomorphism group at the pinned version; the formal proof uses the compact-open topology on $C(Y,Y)$, which for compact metric $Y$ is the topology of the uniform metric.\\
Definition~\ref{def:fat} & \lean{HPC.Effros.NowhereMeagre}, \lean{HPC.Effros.IsFat}.\\
Lemma~\ref{lem:fat-basic} & \lean{HPC.Effros.NowhereMeagre.diff_meagre} (van Mill's Lemma~2.1), \lean{HPC.Effros.NowhereMeagre.closure_subset_closure_diff}, \lean{HPC.Effros.nonempty_interior_closure_of_not_meagre}.\\
The nested construction (Lemma~\ref{lem:core-construction}) & \lean{HPC.Effros.CoreGood}, \lean{HPC.Effros.coreGood_step}, \lean{HPC.Effros.exists_eq_of_core}.\\
Reducing the parameter space (Lemma~\ref{lem:reduced}) & \lean{HPC.Effros.exists_reduced_param} (stated for $P=\N^\N$, which suffices since Mathlib defines analytic sets as continuous images of $\N^\N$).\\
van Mill, Proposition~2.2, relative form (Theorem~\ref{thm:prop22}) & \lean{HPC.Effros.inter_nonempty_of_nowhereMeagre}.\\
van Mill, Proposition~2.2 (Corollary~\ref{cor:fat-meet}) & \lean{HPC.Effros.inter_nonempty_of_isFat}. The hypothesis that the space be nonempty is needed: in the empty space the empty set is analytic and fat (\lean{HPC.Effros.isFat_empty_of_isEmpty}).\\
Lemma~\ref{lem:vm-basis} & \lean{HPC.Effros.exists_vanMill_basis}.\\
van Mill, Lemma~3.1 (Lemma~\ref{lem:vm31}) & \lean{HPC.Effros.vanMill_3_1}.\\
van Mill, Lemma~3.2 (Lemma~\ref{lem:vm32}) & \lean{HPC.Effros.vanMill_3_2}.\\
van Mill, Corollary~3.3 (Lemma~\ref{lem:vm33}) & \lean{HPC.Effros.vanMill_3_3}.\\
van Mill, Lemma~3.4 (Lemma~\ref{lem:vm34}) & \lean{HPC.Effros.vanMill_3_4} (as printed by van Mill, with the closure).\\
van Mill, Proposition~3.5 (Proposition~\ref{prop:vm35}) & \lean{HPC.Effros.vanMill_3_5} (as printed, with the closure).\\
The open mapping principle; Effros (Theorem~\ref{thm:effros-open}) & \lean{HPC.Effros.effros_open_mapping} and \lean{HPC.Effros.isOpenMap_smul_left}.\\
Effros, Ungar (Theorem~\ref{thm:ungar}) & \lean{HPC.Effros.effros_homogeneous}; for a compact subset $M$ of a metric space which is homogeneous in its own topology (the form used for plane continua), \lean{HPC.Effros.effros_homogeneous_set}.\\[4pt]

\end{longtable}


\end{document}
